\documentclass[aos,preprint]{imsart}



% For figures
\usepackage{graphicx} % more modern

\usepackage{amsfonts,amssymb,float}

\DeclareSymbolFont{sfletters}{OML}{cmbrm}{m}{it}

\DeclareMathSymbol{\salpha}{\mathord}{sfletters}{"0B}

\usepackage{mathrsfs}
\RequirePackage{natbib}
\RequirePackage[colorlinks,citecolor=blue,urlcolor=blue]{hyperref}

\usepackage{xr}
\usepackage{xcite}


\externalcitedocument{boot_LLM_aos_galley_SUPP_18july2019}
\externaldocument{boot_LLM_aos_galley_SUPP_18july2019}

\usepackage{mathtools}
%\usepackage{epsfig} % less modern
\usepackage{subfigure} 
% For citations
\RequirePackage{natbib}
\usepackage{comment}
\usepackage{enumerate}
% For algorithms
%\usepackage[named]{algo}
\usepackage{eufrak}

\usepackage{array,varwidth}

%
%
%

\RequirePackage[OT1]{fontenc}
\RequirePackage{amsthm,amsmath}


% settings
%\pubyear{2005}
%\volume{0}
%\issue{0}
%\firstpage{1}
%\lastpage{8}
%\arxiv{arXiv:0000.0000}

\startlocaldefs
\numberwithin{equation}{section}
\theoremstyle{plain}
\newtheorem{theorem}{Theorem}[section]
\newtheorem{lemma}{Lemma}[section]
\newtheorem{proposition}{Proposition}[section]
\newtheorem{assumption}{Assumption}[section]

\newcommand{\M}{\mathcal{M}}
\newcommand{\bt}{\mathbf{t}}
\newcommand{\C}{\mathbf{C}} 
\newcommand{\R}{\mathbb{R}}
\newcommand{\Q}{\mathbb{Q}}
\renewcommand{\hat}[1]{\widehat{#1}}
\renewcommand{\t}[1]{\widetilde{#1}}
\newcommand{\ov}[1]{\overline{#1}}


\newcommand{\I}{\textsc{I}}
\newcommand{\II}{\textsc{II}}
\newcommand{\III}{\textsc{III}}
\newcommand{\bI}{\textup{\texttt{I}}}
\newcommand{\bII}{\textup{\texttt{II}}}
\newcommand{\bIII}{\textup{\texttt{III}}}

%\newcommand{\I}[1]{\mathbbm{1}_{#1}}
\newcommand{\dK}{d_{\textup{K}}}
\newcommand{\dL}{d_{\textup{L}}}
\newcommand{\dLP}{d_{\textup{LP}}}
\newcommand{\dH}{d_{\textup{H}}}
\newcommand{\dTV}{d_{\textup{TV}}}

\newcommand{\Hil}{\mathcal{H}}
\newcommand{\A}{\mathbf{A}}
\newcommand{\J}{\mathcal{J}}
\newcommand{\B}{\mathbf{B}}
\newcommand{\D}{\mathcal{D}}
\newcommand{\G}{\mathbb{G}}
\newcommand{\N}{\mathbb{N}}
\newcommand{\F}{\mathbb{F}}
\renewcommand{\H}{\mathbf{H}}
%\newcommand{\sss}{\S}
%\renewcommand{\SS}{\mathcal{A}}
\newcommand{\FF}{\mathbb{F}}
\newcommand{\X}{\mathbf{X}}
\renewcommand{\S}{\mathcal{A}}
\newcommand{\Pf}{{\bf{Proof: }}}
%\renewcommand{\S}{\mathcal{A}}
\newcommand{\T}{\mathcal{T}}
%\renewcommand{\tsum}{\textstyle\sum}
\newcommand{\ninf}{n \to \infty}
\renewcommand{\P}{\mathbb{P}}
\newcommand{\E}{\mathbb{E}}
\newcommand{\Z}{\mathbf{Z}}
\renewcommand{\O}{\mathcal{O}}
\newcommand{\myRe}{\text{Re}}
\newcommand{\bo}{{\bf{0}}}
\renewcommand{\L}{\mathcal{L}}
\newcommand{\io}{\text{ i.o.}}
\newcommand{\ult}{\text{ ult.}}
\newcommand{\as}{\text{ a.s.}}
\newcommand{\e}{\epsilon}
\newcommand{\ve}{\varepsilon}
\renewcommand{\d}{$\delta$}
\newcommand{\cov}{\text{Cov}}
\newcommand{\tr}{\operatorname{tr}}
\newcommand{\lmin}{\lambda_{\text{min}}}
\newcommand{\lmax}{\lambda_{\text{max}}}
\newcommand{\var}{\operatorname{var}}
\newcommand{\conv}{\operatorname{conv}}
\newcommand{\vol}{\operatorname{vol}}
\newcommand{\hh}{{\bf{H}}}
\newcommand{\av}{{\text{av}}}
\newcommand{\bard}{\overline{\Delta}}
\newcommand{\ttop}{^{\top}}
\newcommand{\invcov}{\text{\emph{invcov}}}
\newcommand{\op}{_{\textup{op}}}
\renewcommand{\dim}{{\text{dim}_*}}
\newcommand{\n}{{\textstyle{\frac{1}{n}}}}
\newcommand{\tsum}{\textstyle\sum}
\newcommand{\ts}{\textstyle}
\newcommand{\1}{{\bf{1}}}
\newcommand{\err}{\text{err}}
\newcommand{\argmin}{\text{argmin}}
\newcommand{\ccirc}{^{\circ}}
\newcommand{\med}{\text{median}}
\newcommand{\opt}{\textup{opt}}
\renewcommand{\dh}{d_{\text{H}}}
\newcommand{\dt}{d_{\text{TV}}}
\newcommand{\Id}{\text{Id}}
\newcommand{\stable}{\text{stable}}
\newcommand{\Log}{\text{Log}}
\newcommand{\tpilot}{\hat{t}_{\text{pilot}}}
\newcommand{\tadapt}{\hat{t}_{\text{adapt}}}
\newcommand{\topt}{\hat{t}_{\text{opt}}}
\newcommand{\bZ}{\mathbf{Z}}
\newcommand{\bW}{\mathbf{W}}
\newcommand{\bX}{\mathbf{X}}
\newcommand{\bSigma}{\mathbf{\Sigma}}
\newcommand{\0}{\mathbf{0}}
\newcommand{\diag}{\textup{diag}}

%%%%%% commands from LIN %%%%%%%%%%%%%%
\newcommand{\hilbert}{\mathcal{H}}
\newcommand{\expect}{\mathbb{E}}
\newcommand{\prob}{\mathrm{Pr}}
\newcommand{\sci}{\hat{\mathcal{I}}}
\newcommand{\innerprod}[2]{\langle#1,#2\rangle}
\newcommand{\siglevel}{\varrho}
\newcommand{\covarop}{\mathcal{C}}
\newcommand{\asympeq}{\asymp}
\newcommand{\transpose}{\textup{T}}
\newcommand{\linrmk}[1]{\textcolor{red}{[LIN: #1]}}
\newcommand{\linrev}[1]{\textcolor{red}{#1}}
\newcommand{\melrmk}[1]{\textcolor{red}{[MEL: #1]}}

%%%%%% END of commands from LIN %%%%%%%

\endlocaldefs
\def\red{\textcolor{red}}
\def\ci{\cite}
\def\cp{\citep}
\bibliographystyle{imsart-nameyear}
\def\references{\bibliography{decay_boot_bib_galley,fda_galley}}


\begin{document}
\newpage
\begin{center}
\Large{Supplementary Material}
\end{center}

\paragraph{Organization of appendices} 
 In Appendix~\ref{app:intro} we prove Proposition~\ref{PROP:DECAY}, and in Appendices~\ref{app:thmg} and~\ref{app:thmboot} we prove Theorems~\ref{THM:G} and~\ref{THM:BOOT} respectively. These proofs rely on numerous technical lemmas, which are stated and proved in Appendix~\ref{app:technical}. Next, the proof of Theorem~\ref{THM:M} for the multinomial model is given in Appendix~\ref{app:mult}, and the associated technical lemmas are given in Appendix~\ref{app:lemmamult}. Lastly, in Appendix~\ref{app:background} we provide statements of background results, and in Appendix~\ref{app:Gaussian} we provide a discussion of related work in the Gaussian approximation literature.

\paragraph{General remarks and notation} Based on the formulation of Theorems~\ref{THM:G},~\ref{THM:BOOT}, and~\ref{THM:M}, it is sufficient to show that these results hold for all large values of $n$, and it will simplify some of the proofs to make use of this reduction. For this reason, it is  understood going forward that $n$ is sufficiently large for any given expression to make sense. Another convention is that all proofs in 
Appendices~\ref{app:thmg},~\ref{app:thmboot},~\ref{app:technical},~\ref{app:mult}, and~\ref{app:lemmamult} will implicitly assume that $p>k_n$ (unless otherwise stated), because once the proofs are given for this case, it will follow that the low-dimensional case where $p\leq k_n$ can be handled as a direct consequence (which is explained on page~\pageref{lowdimcase}).

To fix some notation that will be used throughout the appendices, let $d\in\{1,\dots,p\}$, and define a generalized version of $M$ as
$$M_d=\max_{j\in\J(d)}S_{n,j}/\sigma_j^{\tau_n}.$$
In particular, the statistic $M$ defined in equation~\eqref{eqn:mdef} is the same as $M_p$.
Similarly, the Gaussian and bootstrap versions of $M_d$ are defined as
$$\tilde M_d=\max_{j\in\J(d)}\tilde S_{n,j}/\sigma_j^{\tau_n},$$
and
$$M_d^{\star} =\max_{j\in\J(d)}S_{n,j}^{\star}/\hat\sigma_j^{\tau_n}.$$
 In addition, define the parameter
\begin{equation}
\beta_n=\alpha(1-\tau_n).
\end{equation}
%
Lastly, we will often use the fact that if a random variable $\xi$ satisfies the bound $\|\xi\|_{\psi_1}\leq c$ for some constant $c$ not depending on $n$, then there is another constant $C>0$ not depending on $n$, such that $\|\xi\|_r\leq C\, r$ for all $r\geq 1$~\cite[Proposition 2.7.1]{VershyninHDP}.

\appendix


\section{Proof of Proposition~\ref{PROP:DECAY}} \label{app:intro}
\proof 
It is a standard fact that for any $s\geq 1$, the $\ell_s$ norm dominates its $w\ell_s$ counterpart, and so
$\|\text{diag}(A)\|_{w\ell_s}\leq \|\text{diag}(A)\|_s$.
Next, since $A$ is symmetric, the Schur-Horn Theorem implies that the vector $\text{diag}(A)$ is majorized by $\lambda(A)$~\cite[p.300]{Olkin:2011}. Furthermore, when $s\geq 1$, the function $\|\cdot \|_s$ is Schur-convex on $\R^p$, which means that if  $u\in\R^p$ is majorized by $v\in\R^p$, then $\|u\|_s\leq \|v\|_s$~\cite[p.138]{Olkin:2011}. Hence,
%
$$\|\text{diag}(A)\|_{w\ell_s}\leq \| \lambda(A)\|_s.$$
%
Finally, if $r\in(0,s)$, then for any $v\in\R^p$, the inequality 
$$\|v\|_s \leq \big(\zeta(s/r)\big)^{1/s}\,\|v\|_{w\ell_r}$$ 
holds, where $\zeta(x):=\sum_{j=1}^{\infty}j^{-x}$ for $x>1$. This bound may be derived as in~\citep[p.257]{Johnstone:2017}, 
%
$$\|v\|_s^s \ = \ \sum_{j=1}^p |v|_{(j)}^s \ \leq \ \sum_{j=1}^p \big(\|v\|_{w\ell_r} j^{-1/r}\big)^s \ \leq \ \zeta(s/r)\cdot \|v\|_{w\ell_r}^s,
$$
which completes the proof.
\qed
%
%
%
%
%

\section{Proof of Theorem~\ref{THM:G}}\label{app:thmg}

\proof Consider the inequality
\begin{equation}
d_{\textup{K}}(\mathcal{L}(M_p),\mathcal{L}(\tilde{M}_p)) \, \leq \, \I_n+\II_n+\III_n,
\end{equation}
where we define
\begin{align}
\I_n &= d_{\text{K}}\Big(\mathcal{L}(M_p) \, , \, \mathcal{L}(M_{k_n})\Big)\\[0.2cm]
%
\II_n &=d_{\text{K}}\Big( \mathcal{L}(M_{k_n})\, , \, \mathcal{L}(\tilde M_{k_n})\Big)\\[0.2cm]
%
\III_n &=d_{\text{K}}\Big( \mathcal{L}(\tilde M_{k_n})\, ,\, \mathcal{L}(\tilde M_p)\Big).\label{eqn:IIIdef}
\end{align}
%
Below, we show that the term $\II_n$ is at most of order $n^{-\frac{1}{2}+\delta}$ in Proposition~\ref{prop:bentkus}. Later on, we establish a corresponding result for $\I_n$ and $\III_n$ in Proposition~\ref{prop:IandIII}. Taken together, these results complete the proof of Theorem~\ref{THM:G}.\qed

\begin{proposition}\label{prop:bentkus}
Fix any number $\delta\in(0,1/2)$, and suppose the conditions of Theorem~\ref{THM:G} hold. Then, 
\begin{equation}\label{eqn:IInbound}
\II_n \, \lesssim \,  n^{-\frac 12+\delta}.
\end{equation}
\end{proposition}
%
%
%
\proof 
%
Let $\Pi_{k_n}\in\R^{k_n\times p}$ denote the projection onto the coordinates indexed by $\J(k_n)$. This means that if we write $\J(k_n)=\{j_1,\dots,j_{k_n}\}$ so that $(\sigma_{j_1},\dots,\sigma_{j_{k_n}})=(\sigma_{(1)},\dots,\sigma_{(k_n)})$, then the $l$th row of $\Pi_{k_n}$ is the standard basis vector $e_{j_l}\in\R^p$.
Next, define the diagonal matrix $D_{k_n}=\diag(\sigma_{(1)},\dots,\sigma_{(k_n)})$. It follows that
$$M_{k_n}=\max_{1\leq j\leq k_n} e_j\ttop D_{k_n}^{-\tau_n}\Pi_{k_n} S_n.$$
Define the matrix $\mathfrak{C}\ttop =D_{k_n}^{-\tau_n} \Pi_{k_n}\Sigma^{1/2}$, which is of size $k_n\times p$. Also, let $r$ denote the rank of $\mathfrak{C}$, and note that $r\leq k_n$, since the matrix $\Sigma$ need not be invertible. Next, consider a decomposition
$$\mathfrak{C}=QR,$$
where the columns of $Q\in\R^{p\times r}$ are an orthonormal basis for the image of $\mathfrak{C}$, and $R\in\R^{r\times k_n}$. Hence, if we define the random vector
\begin{equation}\label{eqn:Zrep}
\breve Z=\ts\frac{1}{\sqrt n}\sum_{i=1}^n Q\ttop Z_i,
\end{equation}
then we have
%
$$D_{k_n}^{-\tau_n}\Pi_{k_n} S_n =R\ttop \breve Z.$$
%
It is simple to check that for any fixed $t\in\R$, there exists a Borel convex set $\mathcal{A}_t\subset\R^{r}$ such that $\P(M_{k_n}\leq t)=\P(\breve Z\in\mathcal{A}_t)$. By the same reasoning, we also have $\P(\tilde M_{k_n}\leq t)=\gamma_r(\mathcal{A}_t)$, where $\gamma_{r}$ is the standard Gaussian distribution on $\R^{r}$. Therefore, the quantity $\II_n$  satisfies the bound
%
%
 \begin{equation}
 \begin{split}
 \II_n
 %
 & \ \leq  \ \sup_{\mathcal{A}\in\mathscr{A}}\, \Big| \P\big(\breve Z\in \mathcal{A} \big)- \gamma_{r}(\mathcal{A})\Big|,
 \end{split}
\end{equation}
where $\mathscr{A}$ denotes the collection of all Borel convex subsets of $\R^{r}$. 


We now apply Theorem 1.1 of~\cite{Bentkus:2003} (Lemma~\ref{lem:bentkus}), to handle the supremum above. First observe that
%
the definition of $\breve Z$ in~\eqref{eqn:Zrep} satisfies the conditions of that result, since the terms $Q\ttop Z_1,\dots, Q\ttop Z_n$ are i.i.d.~with zero mean and identity covariance matrix. Therefore,
$$\II_n \ \lesssim \ r^{1/4}\cdot  \E\big[ \|Q\ttop Z_1\|_2^3] \cdot n^{-1/2}.$$
%

It remains to bound the middle factor on the right side. By Lyapunov's inequality, 
%
\begin{equation}\label{eqn:fourthcalc}
\begin{split}
\E\big[\|Q\ttop Z_1\|_2^3\big] &\leq \Big(\E\Big[\big(Z_1\ttop QQ\ttop Z_1\big)^2\Big]\Big)^{3/4}.
%
\end{split}
\end{equation} 
Next, if $v_1,\dots,v_r$ denote the orthonormal columns of $Q$, then we have
$$QQ\ttop = \tsum_{j=1}^r v_jv_j\ttop.$$
Hence, if we put $\zeta_j=Z_1\ttop v_j$, then
\begin{equation}
\begin{split}
 \E\Big[(Z_1\ttop QQ\ttop Z_1)^2\Big] \ &= \ \Big\|\tsum_{j=1}^r \zeta_j^2\Big\|_2^2\\[0.2cm]
 %
 & \ \leq \ \Big(\tsum_{j=1}^r \|\zeta_j^2\|_2\Big)^2\\[0.2cm]
 %
 %
 & \ \lesssim \ k_n^2,
 \end{split}
\end{equation}
where we have used the fact that $r\leq k_n$ and $\|\zeta_j^2\|_2=\|Z_1\ttop v_j\|_4^2\lesssim 1$, based on Assumption~\ref{A:model}.
Combining the last few steps gives $\E[\|Q\ttop Z_1\|_2^3]\lesssim k_n^{6/4}$, and hence
$$ \II_n  \ \lesssim \ k_n^{7/4}n^{-1/2}\ \lesssim \ n^{-\frac 12+\delta},$$
as needed.
\qed 

~\\

\begin{proposition}\label{prop:IandIII}
Fix any number $\delta\in(0,1/2)$, and suppose the conditions of Theorem~\ref{THM:G} hold. Then, 
\begin{equation}\label{IandIII}
\I_n \, \lesssim \, n^{-\frac 12+\delta} \ \ \ \text{ and } \ \ \ \III_n \, \lesssim \, n^{-\frac 12+\delta}.
\end{equation}
\end{proposition}
\proof We only prove the bound for $\I_n$, since the same argument applies to $\III_n$.
%
It is simple to check that for any fixed real number $t$,
$$\Big|\P\Big(\max_{1\leq j\leq p} S_{n,j}/\sigma_j^{\tau_n} \leq t\Big)-\P\Big(\max_{j\in\J(k_n)} S_{n,j}/\sigma_j^{\tau_n} \leq t\Big)\Big| = \ \P\Big(A(t)\cap B(t)\Big),$$
where we define the events
\begin{equation}\label{inclusion}
\small
A(t)=\Big\{\max_{j\in\J(k_n)} S_{n,j}/\sigma_j^{\tau_n} \leq t\Big\} \ \ \ \text{ and } \ \  \ \ B(t)=\Big\{\max_{j\in\J(k_n)^c} S_{n,j}/\sigma_j^{\tau_n}> t\Big\},
\end{equation}
and $\J(k_n)^c$ denotes the complement of $\J(k_n)$ in $\{1,\dots,p\}$.
Also, for any pair of real numbers $t_{1,n}$ and $t_{2,n}$ satisfying $t_{1,n}\leq t_{2,n}$, it is straightforward to check that the following inclusion holds for all $t\in\R$,
\begin{equation}\label{inclusion}
A(t)\cap B(t) \ \subset \ A(t_{2,n})\cup B(t_{1,n}).
\end{equation}
%
Applying a union bound, and then taking the supremum over $t\in\R$, we obtain
$$\I_n\,\leq \, \P(A(t_{2,n}))\,+\, \P(B(t_{1,n})).$$

The remainder of the proof consists in selecting $t_{1,n}$ and $t_{2,n}$ so that $t_{1,n}\leq t_{2,n}$ and that the probabilities $\P(A(t_{2,n}))$ and $\P(B(t_{1,n}))$ are sufficiently small. Below, Lemma~\ref{lem:bounds} shows that if $t_{1,n}$ and $t_{2,n}$ are chosen as 
\begin{align}
t_{1,n}&=c\cdot k_n^{-\beta_n}\cdot \log(n)  \label{eqn:t1}\\[0.3cm]
%
t_{2,n}&=c_{\circ}\cdot \ell_n^{-\beta_n}\cdot\sqrt{\log(\ell_n)},\label{eqn:t2}
\end{align}
for a certain constant $c>0$, and $c_{\circ}$ as in~\eqref{eqn:varcond2}, then  $\P(A(t_{2,n}))$ and $\P(B(t_{1,n}))$ are at most of order $n^{-\frac 12+\delta}$. Furthermore, the inequality $t_{1,n}\leq t_{2,n}$ holds for all large $n$, due to the definitions of $\ell_n$, $k_n$, and $\beta_n$, as well as the condition $(1-\tau_n)\sqrt{\log(n)}\gtrsim 1$.
%
\qed
%
%
%
\begin{lemma}\label{lem:bounds}
Fix any number $\delta\in(0,1/2)$, and suppose the conditions of Theorem~\ref{THM:G} hold. Then, there are positive constants $c$ and $c_{\circ}$, not depending on $n$, that can be selected in the definitions of $t_{1,n}$~\eqref{eqn:t1} and $t_{2,n}$\,\eqref{eqn:t2}, so that
\begin{equation}\label{aboundlem}
\P(A(t_{2,n})) \, \lesssim \, n^{-\frac 12+\delta},\tag{a}
\end{equation}
and
\begin{equation}\label{bboundlem}
\P(B(t_{1,n})) \, \lesssim \,  n^{-1}.\tag{b}
\end{equation}
%
%
\end{lemma}


\paragraph{\textsc{Proof of Lemma~\ref{lem:bounds} part} \eqref{aboundlem}} Due to Proposition~\ref{prop:bentkus} and the fact that $\J(\ell_n)\subset \J(k_n)$, we have
%
\begin{equation}
\begin{split}
\P(A(t_{2,n})) &\ \leq \ \P\Big(\max_{j\in\J(k_n)} \tilde S_{n,j}/\sigma_j^{\tau_n}\leq t_{2,n}\Big)+\II_n\\[0.2cm]
%
& \ \leq \  \ \P\Big(\max_{j\in\J(\ell_n)} \tilde S_{n,j}/\sigma_j^{\tau_n}\leq t_{2,n}\Big)+  c\, n^{-\frac 12+\delta}.
\end{split}
\end{equation}
%
%
%
To bound the probability in the last line, we will make use of the Gaussianity of $\tilde S_n$ to apply certain results based on Slepian's lemma, as contained in Lemmas~\ref{lem:slepian2} and \ref{lem:slepian} below.  As a preparatory step, consider some generic random variables $\{Y_j\}$ and positive scalars $\{a_j\}$ indexed by $\J(\ell_n)$, as well as a constant $b$ such that $\max_{j\in\J(\ell_n)}a_j\leq b$. Then,
\begin{equation}
\begin{split}
\P\Big(\max_{j\in\J(\ell_n)} Y_j\leq t_{2,n}\Big) & \leq \ \P\Big(\max_{j\in\J(\ell_n)}a_jY_j \leq b\, t_{2,n}\Big),
%
\end{split}
\end{equation}
which can be seen by expressing the left side in terms of $\cap_j\{a_jY_j\leq a_j t_{2,n}\}$, and noting that this set is contained in $\cap_j\{a_jY_j\leq b t_{2,n}\}$. Due to Assumption~\ref{A:cor} with $c_{\circ}\in(0,1)$, we have the inequality $\sigma_j^{\tau_n-1}\leq \ell_n^{\beta_n}/c_{\circ}$ for all $j\in\J(\ell_n)$, and so we may apply the previous observation with  $a_j=\sigma_j^{\tau_n-1}$, and $b=\ell_n^{\beta_n}/c_{\circ}$. Furthermore, the definition of $t_{2,n}$ gives $b\, t_{2,n}=\sqrt{\log(\ell_n)}$, and so we if we let $Y_j=\tilde S_{n,j}/\sigma_j^{\tau_n}$, it follows that
$$\P\Big(\max_{j\in\J(\ell_n)} \tilde S_{n,j}/\sigma_j^{\tau_n}\leq t_{2,n}\Big) \ \leq \  \P\Big(\max_{j\in\J(\ell_n)} \tilde S_{n,j}/\sigma_j\leq \sqrt{\log(\ell_n)}\Big).$$
%
The proof is completed by applying the next result (Lemma~\ref{lem:slepian2}) in conjunction with the conditions of Assumption~\ref{A:cor}. (Take $m=\ell_n$ in the statement of Lemma~\ref{lem:slepian2}.)\qed

\paragraph{Remark} The lemma below may be of independent interest, and so we have stated it in a way that can be understood independently of the context of our main assumptions. Also, the constants $1/2$ and $1/3$ in the exponent of the bound~\eqref{eqn:slepian2result} can be improved slightly, but we have left the result in this form for simplicity.

\begin{lemma}\label{lem:slepian2} 
For each integer $m\geq 1$, let $\mathsf{R}=\mathsf{R}(m)$ be a correlation matrix in $\R^{m\times m}$, and let $\mathsf{R}^+=\mathsf{R}^+(m)$ denote the matrix with $(i,j)$ entry given by $\max\{\mathsf{R}_{i,j},0\}$. Suppose the matrix $\mathsf{R}^+$ is positive semi-definite for all $m$, and that there are constants $\e_0\in(0,1)$ and $c>0$, not depending on $m$, such that the inequalities
\begin{align}
\ \max_{i\neq j}\mathsf{R}_{i,j}  & \  \leq \, 1-\e_0\\[0.2cm]
\ts\sum_{i\neq j} \mathsf{R}_{i,j}^+ & \  \leq  \ c\, m\label{eqn:rpluscondn}
\end{align}
hold for all $m$.
Lastly, let $(\zeta_1,\dots,\zeta_m)$ be a Gaussian vector drawn from $N(0,\mathsf{R})$. Then, there is a constant $C>0$, not depending on $m$, such that the inequality
\begin{equation}\label{eqn:slepian2result}
\P\Big(\max_{1\leq j\leq m} \zeta_j\leq \sqrt{\log(m)}\Big) \ \leq \ C\exp\big(-\ts\frac{1}{2}m^{1/3}\big)
\end{equation}
holds for all $m$.
\end{lemma}
\proof It is enough to show that the result holds for all large $m$, because if $m\leq m_0$, then the result is clearly true when $C= \exp(\frac{1}{2}m_0^{1/3})$. To begin the argument,  we may introduce a Gaussian vector $(\xi_1,\dots,\xi_m)\sim N(0,\mathsf{R}^+)$, since the matrix $\mathsf{R}^+$ is positive semi-definite. In turn, the version of Slepian's Lemma given in Lemma~\ref{lem:slepian} leads to
\begin{equation}\label{eqn:splitfactor}
\begin{split}
\P\Big(\max_{1\leq j\leq m}\zeta_j \leq \sqrt{\log(m)}\Big) 
%
& \ \leq \ \P\Big(\max_{1\leq j\leq \ell_n}\xi_j \leq  \sqrt{\log(m)}\Big)\\[0.3cm]
%
 & \ \leq \ K_m\cdot \Phi^{m}\big( \sqrt{\log(m)}\big),
\end{split}
\end{equation}
where we put
$$K_m=  \exp\Bigg\{\sum_{1\leq i<j\leq m} \log\Big(\ts\frac{1}{1-\frac{2}{\pi}\arcsin(\mathsf{R}_{i,j}^+)}\Big)\exp\Big(-\ts\frac{\log(m)}{1+\mathsf{R}_{i,j}^+}\Big)\Bigg\}.$$
Next, we apply the assumption  $\max_{i\neq j}\mathsf{R}_{i,j}^+\leq 1-\e_0$. Since the functions $x\mapsto \log(1/(1-x))$ and  $x\mapsto\frac{2}{\pi}\arcsin(x)$ have bounded derivatives on any closed subinterval of $[0,1)$,  it follows that
 $$\log\Big(\ts\frac{1}{1-\frac{2}{\pi}\arcsin(\mathsf{R}_{i,j}^+)}\Big) \ \leq c\, \mathsf{R}_{i,j}^+,$$
 for some constant $c>0$ not depending on $m$. Therefore, by possibly increasing $c$, the condition~\eqref{eqn:rpluscondn} gives
\begin{equation}
\begin{split}
K_m
%
& \ \leq \ \exp\Big\{ cm \cdot \exp\Big(-\ts\frac{\log(m)}{1+(1-\e_0)}\Big)\Big\}\\[0.3cm]
%
& \ = \ \exp\Big\{cm^{1-\frac{1}{2-\e_0}}\Big\}.
\end{split}
\end{equation}
To bound the earlier factor involving $\Phi^m(\sqrt{\log(m)})$, let $\eta_0\in(0,1)$ be a small constant to be optimized below, and note that the following inequality holds for all sufficiently large $s>0$,
\begin{equation}\label{numericalineq}
\begin{split}
\Phi\Big(\sqrt{(2-\eta_0\e_0)\log(s)}\Big) \ \leq \ 1-\ts\frac{1}{s},
\end{split}
 \end{equation}
%
which may be found in~\cite[p.337]{Massart:2013}. Taking $s=m^{\kappa_0}$ with $\kappa_0:=\frac{1}{2-\eta_0\e_0}$ shows that for all large $m$,
\begin{equation}
\begin{split}
\Phi^{m}\big(\sqrt{\log(m)}\big) & \ \leq  \ \Big(1-\ts\frac{1}{m^{\kappa_0}}\Big)^{m}\\[0.3cm]
%
%
& \ \leq \ \exp\big(-m^{1-\kappa_0}\big).
%
\end{split}
\end{equation}
We now collect the last several steps. If we observe that $\kappa_0<\frac{1}{2-\e_0}$, then the following inequalities hold for all large $m$,
\begin{equation}
\begin{split}
  K_m\cdot \Phi^{m}\big( \sqrt{\log(m)}\big) & \ \leq \ \exp\Big\{cm^{1-\frac{1}{2-\e_0}}-m^{1-\kappa_0}\Big\}\\[0.3cm]
  %
  & \ \leq \ \exp\Big\{-(1-\eta_0)m^{1-\kappa_0}\Big\}.
  %
  \end{split}
\end{equation}
So, by possibly further decreasing $\eta_0$, we have $(1-\kappa_0)>1/3$, as well as $(1-\eta_0)>1/2$. This leads to the stated result.
\qed


~\\

\paragraph{\textsc{Proof of Lemma~\ref{lem:bounds} part}~\eqref{bboundlem}} Define the random variable
$$V=\max_{j\in\J(k_n)^c} S_{n,j}/\sigma_j^{\tau_n},$$
%
and let
$$q=\max\big\{\ts\frac{2}{\beta_n}, \log(n), 3\big\}.$$
Clearly, for any $t>0$, we have the tail bound
\begin{equation}\label{qchebyshev}
\P\big( V\geq t)\leq \frac{\|V\|_q^q}{t^q},
\end{equation}
and furthermore
 \begin{equation}
 \begin{split}
  \|V\|_q^q &= \E\bigg[\Big|\max_{j\in\J(k_n)^c} S_{n,j}/\sigma_j^{\tau_n}\Big|^q\bigg]\\[0.3cm]
  %
  &\leq \sum_{j\in\J(k_n)^c} \sigma_j^{q(1-\tau_n)}\,\E\big[|\ts\frac{1}{\sigma_j}S_{n,j}|^q\big].
  \end{split}
 \end{equation}
By Lemma~\ref{lem:Snjnorm}, we have $\|\frac{1}{\sigma_j}S_{n,j}\|_{q}\leq cq$, and so 
%
  \begin{equation}
  \begin{split}
  %
    \|V\|_q^q \ &\leq \  (c q)^q \sum_{j\in\J(k_n)^c} \sigma_j^{q(1-\tau_n)} \\[0.3cm]
  %
  &\lesssim \ (c q)^q \sum_{j=k_n+1}^p j^{-q \beta_n }\\[0.3cm]
  %
  &\leq \ (cq)^q \int_{k_n}^p x^{-q\beta_n} dx\\[0.3cm]
  %
  &\leq \ \ts\frac{(c q)^q }{q \beta_n-1}\ k_n^{-q \beta_n  +1}, 
 % 
 \end{split}
 \end{equation}
 where we recall $\beta_n=\alpha(1-\tau_n)$, and note that $ q\beta_n\geq 2$, which holds by the definition of $q$.
Hence, if we put $C_n:=\ts\frac{c}{(q\beta_n -1)^{1/q}} \cdot k_n^{1/q}$, then
$$\|V\|_q \leq \ C_n \cdot q\cdot k_n^{-\beta_n}.$$
Furthermore, it is simple to check that $C_n\lesssim 1$, and that the assumption $(1-\tau_n)\sqrt{\log(n)}\gtrsim 1$ implies
$q\lesssim \log(n).$
Therefore, from the inequality~\eqref{qchebyshev} with $t=e\|V\|_q$, as well as the definition of $q$, we obtain
%
%
$$\P\Bigg(V\geq c\cdot \log(n)\cdot k_n^{-\beta_n}\Bigg) \ \leq \ e^{-q} \ \leq \ \ts\frac 1n,$$
for some constant $c>0$ not depending on $n$, as needed.\qed

%

~\\






\section{Proof of Theorem~\ref{THM:BOOT}}\label{app:thmboot} 
\proof Consider the inequality
\begin{equation}
d_{\textup{K}}(\mathcal{L}(\tilde M_p),\mathcal{L}(M_p^{\star}| X)) \ \leq \ \I'_n \ + \ \II'_n(X) \ + \ \III'_n(X),
\end{equation}
where we define
\begin{align}
\I'_n & \ = \  d_{\text{K}}\Big(\mathcal{L}(\tilde M_p) \, , \, \mathcal{L}(\tilde M_{k_n})\Big)\\[0.2cm]
%
\II'_n(X) &\ = \ d_{\text{K}}\Big(\mathcal{L}(\tilde M_{k_n}) \, , \, \mathcal{L}( M^{\star}_{k_n}|X\big)\Big)\\[0.2cm]
%
\III'_n(X) & \ = \ d_{\text{K}}\Big(\mathcal{L}(M^{\star}_{k_n}| X\big) \, , \, \mathcal{L}( M_p^{\star} |X\big)\Big).
\end{align}
Note that $\I_n'$ is deterministic, whereas $\II_n'(X)$ and $\III_n'(X)$ are random variables depending on $X$.
The remainder of the proof consists in showing that each of these terms are at most of order $n^{-\frac 12+\delta}$, with probability at least $1-\frac cn$. The terms $\II_n'(X)$ and $\III_n'(X)$ are handled in Sections~\ref{sec:II'} and~\ref{sec:III'} respectively. The first term $\I_n'$ requires no further work, due to Proposition~\ref{prop:IandIII} (since $\I_n'$ is equal to $\III_n$, defined in equation~\eqref{eqn:IIIdef}).  \qed



%
%
%
%
%
%
%
\subsection{Handling the term $\III_n'(X)$}\label{sec:III'}


The proof of Proposition~\ref{prop:IandIII} can be partially re-used to show that for any fixed realization of $X$, and any real numbers $t_{1,n}'\leq t_{2,n}'$, the following bound holds
\begin{equation}
%
\III_n'(X) \ \leq \ \P\big(A'(t_{2,n}')\big| X\big) \ + \ \P\big(B'(t_{1,n}')\big| X\big),
\end{equation}
where we define the following events for any $t\in\R$,
\begin{equation}\label{inclusionboot}
\small
A'(t)=\Big\{\max_{j\in\J(k_n)} S^{\star}_{n,j}/\hat \sigma_j^{\tau_n}\leq t\Big\} \ \ \ \text{ and } \ \  \ \ B'(t)=\Big\{\max_{j\in\J(k_n)^c} S^{\star}_{n,j}/\hat \sigma_j^{\tau_n}> t\Big\}.
\end{equation} 
Below, Lemma~\ref{lem:Bboundboot} ensures that $t_{1,n}'$ and $t_{2,n}'$ can be chosen so that 
%
the random variables $\P(B'(t_{1,n}')\big|X)$ and $\P(A'(t_{2,n}')\big|X)$ are at most $c n^{-\frac 12+\delta}$, with probability at least $1-\frac cn$. Also, it is straightforward to check that under Assumption~\ref{A:cor}, the choices of $t_{1,n}'$ and  $t_{2,n}'$ given in Lemma~\ref{lem:Bboundboot} satisfy $t_{1,n}'\leq t_{2,n}'$ when $n$ is sufficiently large.


\begin{lemma}\label{lem:Bboundboot}
Fix any number $\delta\in(0,1/2)$, and suppose the conditions of Theorem~\ref{THM:G} hold. Then, there are positive constants $c_1$, $c_2$, and $c$, not depending on $n$,  for which the following statement is true: \\[-0.3cm]

If $t_{1,n}'$ and $t_{2,n}'$ are chosen as
\begin{align}
 t_{1,n}' &= c_1\cdot k_n^{-\beta_n}\cdot\log(n)^{3/2}   \ \ \ \ \text{ and }
 \label{eqn:t1prime}\\[0.2cm]
 %
 t_{2,n}' &=c_2\cdot \ell_n^{-\beta_n}\cdot\sqrt{\log(\ell_n)},
%
\label{eqn:t2prime}
\end{align}
then the events
\begin{equation}\label{abound}\tag{a}
\P(A'(t'_{2,n})\big|X) \leq c\, n^{-\frac 12+\delta}
\end{equation}
and
\begin{equation}\label{bbound}\tag{b}
\P(B'(t'_{1,n})\big| X) \leq n^{-1}
\end{equation}
each hold with probability at least $1-\frac cn$.
%
%
%
\end{lemma}


\paragraph{\textsc{Proof of Lemma~\ref{lem:Bboundboot} part}~\eqref{abound}}Using the definition of $\II_n'(X)$, followed by  $\J(\ell_n)\subset \J(k_n)$, we have
\begin{equation}
\begin{split}
\P(A'(t'_{2,n})|X) 
%
& \ \leq \ \P\Big(\max_{j\in \J(\ell_n)} \tilde S_{n,j}/\sigma_j^{\tau_n}\leq t'_{2,n}\Big) +\II_n'(X).
\end{split}
\end{equation}
Taking $t_{2,n}'=t_{2,n}$ as in~\eqref{eqn:t2}, the proof of Lemma~\ref{lem:bounds} part \eqref{aboundlem} shows that the first term is $\mathcal{O}(n^{-1/2})$. With regard to the second term, Proposition~\ref{prop:IIprime} in the next subsection shows that  there is a constant $c>0$ not depending on $n$ such that the event
\begin{equation}
\II_n'(X) \leq c\, n^{-\frac 12+\delta}
\end{equation}
holds with probability at least $1-\frac cn$. This completes the proof.\qed









%%%%%%%%%%%%%%%%%%%%%%%%%%



~\\


\paragraph{\textsc{Proof of Lemma~\ref{lem:Bboundboot} part}~\eqref{bbound}} Define the random variable
\begin{equation}\label{eqn:Vstar}
V^{\star}:=\max_{j\in\J(k_n)^c} S_{n,j}^{\star}/\hat\sigma_j^{\tau_n},
\end{equation}
and as in the proof of Lemma~\ref{lem:bounds}\eqref{bboundlem}, let 
$q =\max\big\{\ts\frac{2}{\beta_n},\log(n),3\big\}.$
The idea of the proof is to construct a function $b(\cdot)$ such that the following bound holds for every realization of $X$,
$$\Big( \E\big[|V^{\star}|^q\big | X\big]\Big)^{1/q} \leq  b(X),$$
and then Chebyshev's inequality gives the following inequality for any number $b_n$ satisfying $b(X)\leq b_n$,
$$\P\Big(V^{\star} \geq e b_n\, \Big | X\Big)  \ \leq \ e^{-q} \ \leq \ \ts\frac 1n.$$
%
In turn, we will derive an expression for $b_n$ such that the event $\{b(X)\leq b_n\}$ holds with high probability.
This will lead to the statement of the lemma, because it will turn out that $t_{1,n}'\asymp b_n$.\\

To construct the function $b(\cdot)$, observe that the initial portion of the proof of Lemma~\ref{lem:bounds}\eqref{bboundlem} shows that for any realization of $X$,
\begin{equation}\label{eqn:Vstarbound}
 \E\big[|V^{\star}|^q\big | X\big]\ \leq   \sum_{j\in\J(k_n)^c} \hat\sigma_j^{q(1-\tau_n)}\,\E\big[|\ts\frac{1}{\hat \sigma_j}S_{n,j}^{\star}|^q|X\big].
\end{equation}
%
Next, Lemma~\ref{lem:Snjnorm} ensures that for every $j\in\{1,\dots,p\}$, the event
%
$$ \E\big[|\ts\frac{1}{\hat \sigma_j}S_{n,j}^{\star}|^q|X\big] \ \leq (c\,q)^q, $$
holds with probability 1.
Consequently, if we let $s=q(1-\tau_n)$ and consider the random variable
\begin{equation}
\hat{\mathfrak{s}}:=\Bigg(\sum_{j\in\J(k_n)^c}\hat\sigma_j^{s}\Bigg)^{\frac{1}{s}},
\end{equation}
as well as
$$b(X):=c\cdot q\cdot  \hat{\mathfrak{s}}^{(1-\tau_n)},$$
then we obtain the bound
\begin{equation}
\Big( \E\big[|V^{\star}|^q\big | X\big]\Big)^{1/q} \ \leq b(X),
\end{equation}
with probability 1.
%
To proceed, Lemma~\ref{lem:mathfrak} implies 
\begin{equation}
\P\bigg(b(X) \geq  q\cdot \ts\frac{(c\sqrt q)^{1-\tau_n}}{(q\beta_n-1)^{1/q}}\cdot k_n^{-\beta_n+1/q}\bigg) \ \leq e^{-q}\leq \frac 1n,
\end{equation}
for some constant $c>0$ not depending on $n$.
%
By weakening this tail bound slightly, it can be simplified to
%
\begin{equation}
\P\bigg(b(X) \geq C_n'\cdot  q^{3/2} \cdot k_n^{-\beta_n}\bigg) \ \leq \ts\frac 1n,
\end{equation}
where  $C_n':=\ts\frac{ c \,k_n^{1/q}}{(q\beta_n-1)^{1/q}}$, and we recall $\beta_n=\alpha(1-\tau_n)$.
To simplify further, it can be checked that $C_n'\lesssim 1$, and that the assumption $(1-\tau_n)\sqrt{\log(n)}\gtrsim 1$ gives $q\lesssim \log(n)$. It follows that there is a constant $c$ not depending on $n$ such that if
$$b_n:=c\cdot \log(n)^{3/2} \cdot  k_n^{-\beta_n},$$
then
\begin{equation}
\P(b(X)\geq  b_n)  \ \leq \ts\frac 1n,
\end{equation}
which completes the proof.\qed



\subsection{Handling the term $\II_n'(X)$}\label{sec:II'}
\begin{proposition}\label{prop:IIprime}
Fix any number $\delta\in(0,1/2)$, and suppose the conditions of Theorem~\ref{THM:G} hold. Then, there is a constant $c>0$ not depending on $n$ such that the event
\begin{equation}
\II_n'(X) \leq c\, n^{-\frac 12+\delta}
\end{equation}
holds with probability at least $1-\frac cn$.
%
\end{proposition}

\proof Define the random variable
\begin{equation}
\breve M_{k_n}^{\star}:=\max_{j\in \J(k_n)} S_{n,j}^{\star}/\sigma_j^{\tau_n},
\end{equation}
which differs from $M_{k_n}^{\star}$, since $\sigma_j^{\tau_n}$ is used in place of $\hat\sigma_j^{\tau_n}$.
 Consider the triangle inequality
 \begin{equation}\label{eqn:II'}
 \small
\II_n'(X) \ \leq \ \dK\Big(\L(\tilde M_{k_n})\, , \L(\breve M_{k_n}^{\star}|X)\Big) \ + \ \dK \Big(\L(\breve M_{k_n}^{\star}|X) \, , \, \L(M_{k_n}^{\star}|X)\Big).
\end{equation}
%
The two terms on the right will bounded separately. 

To address the first term on the right side of~\eqref{eqn:II'}, we will apply Lemma~\ref{lem:hellinger}, for which a substantial amount of notation is needed. Recall the matrix $\mathfrak{C}=\Sigma^{1/2}\Pi_{k_n}\ttop D_{k_n}^{-\tau_n}$ of size $p\times k_n$, where the projection matrix $\Pi_{k_n}\in\R^{k_n\times p}$ is defined in the proof of Proposition~\ref{prop:bentkus}.  Note that $\tilde M_{k_n}$ is the coordinate-wise maximum of a Gaussian vector drawn from $N(0,\mathfrak{S})$, with $\mathfrak{S}=\mathfrak{C}\ttop\mathfrak{C}$. To address $\breve{M}_{k_n}^{\star}$, let
$$W_n=\ts\frac 1n \sum_{i=1}^n  (Z_i-\bar Z)(Z_i-\bar Z)\ttop$$ 
where $\bar Z=\frac 1n \sum_{i=1}^n Z_i$, and observe that $\breve{M}_{k_n}^{\star}$ is the coordinate-wise maximum of Gaussian vector drawn from $N(0,\breve{\mathfrak{S}})$, with $\breve{\mathfrak{S}}=\mathfrak{C}\ttop W_n\mathfrak{C}$. Next, consider the s.v.d.,
$$\mathfrak{C}=U\Lambda V\ttop,$$
where if $r$ denotes the rank of $\mathfrak{C}$, then we may take $U\in\R^{p\times r}$ to have orthonormal columns, $\Lambda\in\R^{r\times r}$ to be invertible, and $V\ttop\in\R^{r\times k_n}$ to have orthonormal rows. In order to apply Lemma~\ref{lem:hellinger} for a given realization of $\breve{\mathfrak{S}}$, it is necessary that the columns of $\mathfrak{S}$ and $\breve{\mathfrak{S}}$ span the same subspace of $\R^{k_n}$. This occurs with probability at least $1-\frac cn$, because the matrix $\breve{\mathfrak{S}}$ is equal to $V\Lambda (U\ttop W_n U)\Lambda V\ttop$, and the matrix $(U\ttop W_n U)$ is invertible with probability at least $1-\frac cn$ (due to Lemma~\ref{lem:white}). Another ingredient for applying Lemma~\ref{lem:hellinger} is the following algebraic relation, which is a direct consequence of the definitions just introduced,
\begin{equation}
\Big(V\ttop \mathfrak{S} V\Big)^{-1/2}\Big( V\ttop \breve{\mathfrak{S}}V\Big) \Big(V\ttop\mathfrak{S}V\Big)^{-1/2} \ = \ U\ttop W_n U.
\end{equation}
Building on this relation, Lemma~\ref{lem:hellinger} shows that if the event
 \begin{equation}\label{eqn:tempevent2}
 \|U\ttop W_n U-\mathbf{I}_{r}\|_{\text{op}} \ \leq \ \e,
 \end{equation}
 holds for some number $\e>0$, then the event
\begin{equation}
 \dK\Big(\L(\tilde M_{k_n})\, , \L(\breve M_{k_n}^{\star}|X)\Big) \leq c\cdot  k_n^{1/2} \cdot \e
\end{equation}
also holds, where $c>0$ is a constant not depending on $n$ or $\e$. Thus, it remains to specify $\e$ in the event~\eqref{eqn:tempevent2}. For this purpose, Lemma~\ref{lem:white} shows that if $\e=c \cdot n^{-1/2}\cdot k_n\cdot \log(n)$, then the event~\eqref{eqn:tempevent2} holds with probability at least $1-\frac cn$. 
%
So, given that 
$$n^{-1/2}\cdot k_n^{3/2}\cdot \log(n) \ \lesssim \ n^{-\frac 12+\delta},$$ 
the first term in the bound~\eqref{eqn:II'} requires no further consideration.\\

To deal with the second term in~\eqref{eqn:II'}, we proceed by considering the general inequality
\begin{equation}
\dK(\L(\xi),\L(\zeta)) \ \leq \ \sup_{t\in\R}\P\big(|\zeta-t|\leq r\big)  \ + \  \P(|\xi-\zeta|>r),
\end{equation}
which holds for any random variables $\xi$ and $\zeta$, and any real number $r>0$ (cf.~\citet[Lemma 2.1]{CCK:SPA}). Specifically, we will let $\L(\breve M_{k_n}^{\star}|X)$ play the role of $\L(\xi)$, and let $\L(M_{k_n}^{\star}|X)$ play the role of $\L(\zeta)$. In other words, we need to establish an anti-concentration inequality for $\mathcal{L}(M_{k_n}^{\star}|X)$, as well as a coupling inequality for $ M_{k_n}^{\star}$ and $\breve M_{k_n}^{\star}$, conditionally on $X$.

To establish the coupling inequality, if we put
\begin{equation}
r_n= c \cdot n^{-1/2}\cdot \log(n)^{5/2},
\end{equation}
for a suitable constant $c$ not depending on $n$, then Lemma~\ref{lem:coupling} shows that the event
\begin{equation}\label{eqn:maincouple}
\P\Big(\big|\breve M_{k_n}^{\star} - M_{k_n}^{\star}\big| > r_n\, \Big| X\Big)  \ \leq  \ \ts\frac cn
\end{equation}
holds with probability at least $1-\frac cn$.


Lastly, the anti-concentration inequality can be derived from Nazarov's inequality (Lemma~\ref{lem:Nazarov}), since $M_{k_n}^{\star}$ is obtained from a Gaussian vector, conditionally on $X$. For this purpose, let 
\begin{equation}
\hat{\underline{\sigma}}_{k_n}=\min_{j\in\mathcal{J}(k_n)}\hat\sigma_j.
\end{equation}
In turn, Nazarov's inequality implies that the event
\begin{equation}
\sup_{t\in\R}\, \P\Big(|M_{k_n}^{\star} -t|\leq r_n \, \Big| X\Big) \ \leq \ c\cdot \frac{r_n}{\hat{\underline{\sigma}}_{k_n}^{1-\tau_n}}\cdot \sqrt{\log(k_n)},
\end{equation}
holds with probability 1.
Meanwhile, Lemma~\ref{lem:cor} and Assumption~\ref{A:cor} imply that the event
\begin{equation}\label{eqn:minsigevent}
\frac{1}{\hat{\underline{\sigma}}_{k_n}^{1-\tau_n}} \ \leq c\, k_n^{\beta_n}
\end{equation}
holds with probability at least $1-\frac cn$. 
%
Combining the last few steps, we conclude that the following bound holds with probability at least $1-\frac cn$,
\begin{equation}
\small
\begin{split}
\sup_{t\in\R}\, \P\Big(|M_{k_n}^{\star} -t|\leq r_n \, \Big| X\Big)  \  & \leq \ts c\cdot n^{-1/2}\cdot k_n^{\beta_n}\cdot \log(n)^{5/2}\cdot\sqrt{\log(k_n)}\\[0.2cm]
%
& \ \leq  c\,n^{-\frac 12 +\delta},
\end{split}
\end{equation}
as needed.
%
\qed


\section{Technical lemmas for Theorems~\ref{THM:G} and~\ref{THM:BOOT}}\label{app:technical}
\begin{lemma}\label{lem:sighatrnorm}
Fix any number $\delta\in(0,1/2)$, and suppose the conditions of Theorem~\ref{THM:G} hold. Also, let $q=\max\{\frac{2}{\beta_n}, \log(n),3\}$. Then, there is a constant $c>0$ not depending on $n$, such that for any $j\in\{1,\dots,p\}$, we have
\begin{equation}
\|\hat\sigma_{j}\|_q \ \leq \ c \cdot \sigma_j \cdot \sqrt{q}.
\end{equation}
\end{lemma}
\proof Define the vector $u:=\frac{1}{\sigma_j}\Sigma^{1/2}e_j\in\R^p$, which satisfies $\|u\|_2=1$. Observe that
\begin{equation}\label{eqn:sigsteps}
\begin{split}
\ts\frac{1}{\sigma_j}\|\hat\sigma_{j}\|_q &= \ \Bigg\|\Big(\ts\frac 1n \sum_{i=1}^n (Z_i\ttop u)^2-(\bar Z\ttop u)^2\Big)^{1/2}\Bigg\|_q\\[0.2cm]
%
&\leq \ \bigg\|\Big(\ts\frac 1n \sum_{i=1}^n ( Z_i\ttop u)^2\Big)^{1/2}\bigg\|_q\\[0.2cm]%
&= \ \Big\|\ts\frac 1n \sum_{i=1}^n (Z_i\ttop u)^2\Big\|_{q/2}^{1/2}\ .
\end{split}
\end{equation}
Since the random variables $(Z_1\ttop u)^2,\dots, (Z_n\ttop u)^2$ are independent~and non-negative, part \emph{(i)} of Rosenthal's inequality in Lemma~\ref{lem:rosenthal} implies the $L^{q/2}$ norm in the last line satisfies
\begin{equation}
\small
 \Big\|\ts\frac 1n \sum_{i=1}^n (Z_i\ttop u)^2\Big\|_{q/2} \ \leq \ c \cdot q \cdot \max\bigg\{\big\|(Z_1\ttop u)^2\big\|_1\, , \, n^{-1+2/q} \big\| (Z_1\ttop u)^2\big\|_{q/2}\bigg\},
\end{equation}
for an absolute constant $c>0$.
For the first term inside the maximum, observe that since $\|u\|_2=1$ and $Z_1$ is isotropic, we have $\|(Z_1\ttop u)^2\|_1=1$. To handle the second term inside the maximum, Assumption~\ref{A:model} implies $\|(Z_1\ttop u)^2\|_{q/2}\lesssim q^2$.
Combining the last few steps, and noticing the square root on the $L^{q/2}$ norm in the last line of~\eqref{eqn:sigsteps}, we obtain 
\begin{equation}
\begin{split}
\ts\frac{1}{\sigma_j}\|\hat\sigma_{j}\|_q & \ \lesssim \sqrt q \cdot \max\Big\{ 1\, , \, n^{-1/2+1/q}q\Big\},
\end{split}
\end{equation}
and this implies the statement of the lemma.\qed



\begin{lemma}\label{lem:mathfrak}
Fix any number $\delta\in(0,1/2)$, and suppose the conditions of Theorem~\ref{THM:G} hold. Also, let $q=\max\{\frac{2}{\beta_n},\log(n),3\}$, and $s=q(1-\tau_n)$, and consider the random variables $\hat{\mathfrak{s}}$ and $\hat{\mathfrak{t}}$ defined by
$$\hat{\mathfrak{s}}=\Bigg(\sum_{j\in\J(k_n)^c} \hat\sigma_{j}^s\Bigg)^{1/s} \text{ \ \ \ \  and \ \ \ \ } \hat{\mathfrak{t}}=\Bigg(\sum_{j\in\J(k_n)} \hat\sigma_{j}^s\Bigg)^{1/s}.$$
Then,
there is a constant $c>0$ not depending on $n$ such that 
%
\begin{equation}
\P\Bigg(\hat{\mathfrak{s}}\geq \ts\frac{c\sqrt q}{(q\beta_n-1)^{1/s}}\cdot k_n^{-\alpha+1/s}\Bigg) \ \leq \  e^{-q},
\end{equation}
and
\begin{equation}\label{eqn:boundmathfrakt}
\P\Big(\, \hat{\mathfrak{t}}\, \geq \ts\frac{c \sqrt q}{(q\beta_n-1)^{1/s}}\Big) \ \leq \  e^{-q}.
\end{equation}
\end{lemma}
\proof 
%
In light of the Chebyshev inequality $\P\big(\hat{\mathfrak{s}}\geq e\|\hat{\mathfrak{s}}\|_q\big) \leq e^{-q}$, it suffices to bound $\|\hat{\mathfrak{s}}\|_q$ (and similarly for $\hat{\mathfrak{t}}$). We proceed by direct calculation,
\begin{equation*}
\footnotesize
\begin{split}
\|\hat{\mathfrak{s}}\|_q \ & = \ \Bigg\|\sum_{j\in\J(k_n)^c}\hat\sigma_{j}^s\Bigg\|_{q/s}^{1/s}\\[0.2cm]
%
&\leq \ \Bigg(\sum_{j\in\J(k_n)^c} \big\|\hat\sigma_{j}^s\big\|_{q/s}\Bigg)^{1/s} \ \ \ \ \ \ \ \text{(triangle inequality for $\|\cdot\|_{q/s}$, with $q/s\geq 1$)}\\[0.2cm]
%
&=\ \ \Bigg(\sum_{j\in\J(k_n)^c} \big\|\hat\sigma_{j}\big\|_q^s\Bigg)^{1/s}\\[0.2cm]
%
&\lesssim \ \sqrt{q} \cdot\Bigg( \sum_{j\in\J(k_n)^c} \sigma_j^{s}\Bigg)^{1/s} \ \ \ \ \ \ \ \ \ \text{(Lemma~\ref{lem:sighatrnorm})}\\[0.2cm]
%
%
&\lesssim \ \sqrt{q} \cdot \Bigg(\int_{k_n}^p x^{-s\alpha}dx\Bigg)^{1/s}  \\[0.2cm]
%
%
%
&\lesssim \ \sqrt{q}\cdot  \frac{k_n^{-\alpha+1/s}}{(s\alpha-1)^{1/s}},
%
\end{split}
\end{equation*}
\normalsize
and in the last step we have used the fact that $s\alpha=q\beta_n>1$, which holds since $q$ is defined to satisfy $q\beta_n>1$. The calculation for $\hat{\mathfrak{t}}$ is essentially the same, except that we use $\sum_{j\in\J(k_n)}\sigma_j^s\lesssim 1$.
\qed
~\\


\paragraph{Remark} The following result is a variant of Lemma A.7 in the paper~\citep{Zhilova:2015}.
\begin{lemma}\label{lem:hellinger}
Let $A$ and $B$ be positive semi-definite matrices in \smash{$\R^{d\times d}\setminus\{0\}$} whose columns span the same subspace of $\R^d$.
Define two multivariate normal random vectors $\xi\sim N(0, A)$ and $\zeta\sim N(0,B)$. 
Let $r\leq d$ be the dimension of the subspace spanned by the columns of $A$ and $B$, and let $Q\in\R^{d\times r}$ have columns that are an orthonormal basis for this subspace. Define the $r\times r$ positive definite matrices $\tilde A=Q\ttop A Q$ and $\tilde B=Q\ttop B Q$, and let $H$ be any square matrix satisfying $H\ttop H=\tilde A$. Finally, let $\e>0$ be a number such that $\|(H^{-1})\ttop \tilde B (H^{-1}) -\mathbf{I}_r\|\op \leq \e$. Then, there is an absolute constant $c>0$ such that
%
\begin{equation}\label{hellenger}
 \sup_{t\in\R}\bigg|\P\Big(\max_{1\leq j\leq d} \xi_j\leq t\Big)-\P\Big(\max_{1\leq j\leq d} \zeta_j\leq t\Big)\bigg| \ \leq  c\sqrt{r}\, \e.
\end{equation}
\end{lemma}


\proof 
We may assume that $\sqrt{r} \,\e\leq 1/2$, for otherwise the claim trivially holds with $c=2$. Define the $r$-dimensional random vectors $\tilde \xi=Q\ttop  \xi$ and $\tilde \zeta=Q\ttop \zeta$. As a consequence of the assumptions, the random vector $\xi$ lies in the column-span of $Q$ almost surely, which gives $Q\tilde \xi=\xi$ almost surely. It follows that for any $t\in\R$, the event $\{\max_{1\leq j\leq d} \xi_j \leq t\}$ can be expressed as $\{\tilde \xi\in\mathcal{A}_t\}$ for some Borel set $\mathcal{A}_t\subset \R^r$. Likewise, we also have  $\{\max_{1\leq j\leq d} \zeta_j \leq t\} = \{\tilde \zeta\in\mathcal{A}_t\}$. Hence, the left hand side of~\eqref{hellenger} is upper-bounded by the total variation distance between $\L(\tilde \xi)$ and $\L(\tilde \zeta)$, and in turn, Pinsker's inequality implies this is upper-bounded by $c\sqrt{d_{\textup{KL}}(\L(\tilde \zeta),\L(\tilde \xi))}$, where $c>0$ is an absolute constant, and $d_{\textup{KL}}$ denotes the KL divergence. Since the random vectors $\tilde \xi$ and $\tilde \zeta$ are Gaussian, the following exact formula is available if we let $\tilde C=(H\ttop)^{-1}\tilde B(H^{-1})-\mathbf{I}_r$, 
\begin{equation}
\begin{split}
 d_{\textup{KL}}(\L(\tilde \zeta),\L(\tilde \xi)) & \ = \ \ts\frac{1}{2}\Big(\tr(\tilde C)-\log\det(\tilde C+\mathbf{I}_r)\Big)\\[0.2cm]
 %
 & \ = \ \ts\frac{1}{2}\sum_{j=1}^r \lambda_j(\tilde C)-\log(\lambda_j(\tilde C)+1).
 \end{split}
\end{equation}
Using the basic inequality  $|x-\log(x+1)|\leq x^2/(1+x)$ that holds for any $x\in(-1,\infty)$, as well as the condition $|\lambda_j(\tilde C)|\leq \e\leq 1/2$, we have
\begin{equation}
\begin{split}
 d_{\textup{KL}}(\L(\tilde \zeta),\L(\tilde \xi)) & \ \leq \ c \, r\, \|\tilde C\|\op^2\\[0.2cm]
 %
 & \ \leq c \, r \, \e^2,
 \end{split}
\end{equation}
for some absolute constant $c>0$. \qed






\begin{lemma}\label{lem:Snjnorm}
Suppose the conditions of Theorem~\ref{THM:G} hold, and let $q=\max\{\frac{2}{\beta_n}, \log(n),3\}$.
 Then, there is a constant $c>0$ not depending on $n$ such that for any $j\in\{1,\dots,p\}$, we have
\begin{equation}\label{eqn:firstsnjnorm}
\|\ts\frac{1}{\sigma_j}S_{n,j}\|_q \leq c\, q,
\end{equation}
and the following event holds with probability 1,
\begin{equation}
\Big(\E\big[|\ts\frac{1}{\hat \sigma_j}S_{n,j}^{\star}|^q|X\big] \Big)^{1/q} \ \leq c\, q.
\end{equation}
\end{lemma}
\proof We only prove the first bound, since the second one can be obtained by repeating the same argument, conditionally on $X$.
%
Since $q>2$, Lemma~\ref{lem:rosenthal} gives
\begin{equation}\label{eqn:lemrosenthalfirstD}
 \|\ts\frac{1}{\sigma_j}S_{n,j}\|_q \ \lesssim \ q\cdot \max\Big\{  \|\ts\frac{1}{\sigma_j}S_{n,j}\|_2 \, , \,  n^{-1/2+1/q} \|\ts\frac{1}{\sigma_j}(X_{1,j}-\mu_j)\|_q\Big\}.
\end{equation}
Clearly,
\begin{equation}
 \|\ts\frac{1}{\sigma_j}S_{n,j}\|_2^2 \ = \  \var(\ts\frac{1}{\sigma_j}S_{n,j}) \ = \ 1.
\end{equation}
Furthermore, if we define the vector $u:=\ts\frac{1}{\sigma_j}\Sigma^{1/2}e_j$ in $\R^p$, which satisfies $\|u\|_2=1$, then 
\begin{equation}
\begin{split}
\big\|\ts\frac{1}{\sigma_j}(X_{1,j}-\mu_j)\big\|_q
%
%
&=\big\| Z_1\ttop u\big\|_q \ \lesssim  \ q \\[0.2cm]
%
%
\end{split}
\end{equation}
where the last step  follows from Assumption~\ref{A:model}. Applying the work above to the bound~\eqref{eqn:lemrosenthalfirstD} gives
\begin{equation}
 \|\ts\frac{1}{\sigma_j}S_{n,j}\|_q \ \lesssim \ q\cdot \max\Big\{ 1, n^{-1/2+1/q}\cdot q\Big\}.
\end{equation}
Finally, the stated choice of $q$ implies that the right side in the last display is of order $q$.\qed



\begin{lemma}\label{lem:white}
 Let the random vectors $Z_1,\dots,Z_n\in\R^p$ be as in Assumption~\ref{A:model}, and let $Q\in \R^{p\times r}$ be a fixed matrix having orthonormal columns with $r\leq k_n$. Lastly, let
 \begin{equation}\label{eqn:Wndef}
 W_n=\frac 1n \sum_{i=1}^n  (Z_i-\bar Z)(Z_i-\bar Z)\ttop,
 \end{equation}
 where $\bar Z=\frac 1n \sum_{i=1}^n Z_i$. Then, there is a constant $c>0$ not depending on $n$, such that the event
  \begin{equation}\label{eqn:opnormevent}
 \big\|Q\ttop W_nQ-\mathbf{I}_{r}\big\|\op \ \leq \ \ts\frac{c\log(n)k_n}{\sqrt n},
 \end{equation}
 holds with probability at least $1-\frac cn$. 
\end{lemma}
\proof 
Let $\e\in (0,1/2)$, and let $\mathcal{N}$ be an $\e$-net (with respect to the $\ell_2$-norm) for the unit $\ell_2$-sphere in $\R^{r}$. It is well known that $\mathcal{N}$ can be chosen so that $\text{card}(\mathcal{N})\leq (3/\e)^{r}$, and the inequality
$$\big\|Q\ttop W_n Q- \mathbf{I}_{r}\big\|\op \leq \ts\frac{1}{1-2\e}\cdot \displaystyle \max_{u\in\mathcal{N}} \Big|u\ttop\Big(Q\ttop W_n Q-\mathbf{I}_{r}\Big)u\Big|,$$
holds with probability 1~\citep[Lemmas 5.2 and 5.4]{Vershynin:2012}. For a fixed $u\in\mathcal{N}$, put $\xi_{u,i}:=Z_i\ttop Qu$,
and consider the simple algebraic relation
\begin{equation}
u\ttop\Big(Q\ttop W_n Q-\mathbf{I}_{r}\Big)u = \underbrace{\Big(\ts\frac{1}{n}\sum_{i=1}^n \xi_{i,u}^2-1\Big)}_{=: \Delta(u)}- \underbrace{\Big(\ts\frac{1}{n}\sum_{i=1}^n \xi_{i,u}\Big)^2}_{=:\Delta'(u)}.
\end{equation}
%
We will show that both terms on the right side are small with high probability, and then take a union bound over $u\in\mathcal{N}$. The high-probability bounds will be obtained by using Lemma~\ref{lem:rosenthal} to control $\|\Delta(u)\|_q$ and $\|\Delta'(u)\|_q$ when $q$ is sufficiently large.\\

To apply Lemma~\ref{lem:rosenthal}, first observe the following bounds, which are consequences of Assumption~\ref{A:model},
\begin{equation}
\|\xi_{i,u}\|_q \ \lesssim \  q, 
\end{equation}
and
\begin{equation}
\begin{split}
\|\xi_{i,u}^2-1\|_q \ \lesssim \ q^2.
\end{split} 
\end{equation}
Therefore, when $q>2$, Lemma~\ref{lem:rosenthal} gives
\begin{equation}\label{eqn:justbefore}
\begin{split}
\|\Delta(u)\|_q  & \ \lesssim \ q \max\Big\{ \|\Delta(u)\|_2 \, , \, \ts\frac{1}{n}\big(\tsum_{i=1}^n \|\xi_{i,u}^2-1\|_q^q\big)^{1/q}\Big\}\\[0.3cm]
%
& \lesssim \ q \max\Big\{ \ts\frac{1}{\sqrt n} \, , \, n^{-1+1/q} \cdot q^2\Big\}.
\end{split}
\end{equation} 
Due to Chebyshev's inequality,
$$\P\Big( |\Delta(u)| \ \geq e \|\Delta(u)\|_q\Big) \ \leq \ e^{-q},$$
and so if we take $q= \max\{C\log(n)k_n,3\}$ for some constant $C>0$ to be tuned below, then~\eqref{eqn:justbefore} gives $\|\Delta(u)\|_q\lesssim q/\sqrt{n}$, and the following inequality holds for all $u\in\mathcal{N}$,
$$\P\Big( |\Delta(u)| \ \geq \ts\frac{c\log(n) k_n}{\sqrt n}\Big) \ \leq \ \exp\Big\{-C\log(n) k_n\Big\}.$$
%
The random variable $\Delta'(u)$ can be analyzed with a similar set of steps, which leads to the following inequality for all $u\in\mathcal{N}$,
\begin{equation}
\P\bigg( |\Delta'(u)| \ \geq c\Big(\ts\frac{\log(n) k_n}{\sqrt n}\Big)^2\bigg) \ \leq \ \exp\Big\{-C\log(n) k_n\Big\}.
\end{equation}
%
Combining the previous work with a union bound, if we consider the choice $\e=\min\{c\log(n)k_n/\sqrt{n},\ts\frac{1}{4}\}$, then 
$$ \P\Big(\big\|Q\ttop W_n Q- \mathbf{I}_{r}\big\|\op \geq \e\Big) \ \leq \ 2\exp\Big\{ - C\cdot k_n\cdot\log(n) \ + \ r\cdot \log(3/\e)\Big\}.$$
%
Finally, choosing $C$ sufficiently large implies the stated result.\qed
 
\paragraph{Remark} For the next results, define the correlation 
$$\rho_{j,j'}=\ts\frac{\Sigma_{j,j'}}{\sigma_j\sigma_{j'}},$$ 
and its sample version 
$$\hat{\rho}_{j,j'}=\ts\frac{\hat\Sigma_{j,j'}}{\hat\sigma_j\hat\sigma_{j'}},$$
 for any $j,j'\in\{1,\dots,p\}$.
\begin{lemma}\label{lem:cor}
Suppose the conditions of Theorem~\ref{THM:G} hold. Then, there is a constant $c>0$ not depending on $n$ such that the three events
\begin{equation}\label{eqn:firstcor}
 \max_{j\in\J(k_n)}\Big| \ts\frac{\hat\sigma_j}{\sigma_j}-1\Big| \leq \ts \frac{c\log(n)}{\sqrt n},
\end{equation}
%
\begin{equation}
 \min_{j\in\J(k_n)}\hat\sigma_j^{1-\tau_n} \ \geq \ \Big(\min_{j\in\J(k_n)} \sigma_j^{1-\tau_n}\Big)\cdot \Big(1-\ts\frac{c\log(n)}{\sqrt n}\Big),
\end{equation}
and
\begin{equation}\label{eqn:firstcor}
 \max_{j,j'\in\J(k_n)}\big| \hat \rho_{jj'}-\rho_{jj'}\big| \leq \ts \frac{c\log(n)}{\sqrt n}
\end{equation}
each
hold with probability at least $1-\ts\frac{c}{n}$.
\end{lemma}
\proof The result is a direct consequence of Lemma~\ref{lem:corbasic} below. The details are essentially algebraic manipulations, and so are omitted.\qed



\begin{lemma}\label{lem:corbasic}
Suppose the conditions of Theorem~\ref{THM:G} hold, and fix any two (possibly equal) indices $j,j'\in\{1,\dots,p\}$.
 Then, for any number $\kappa\geq 1$, there are positive constants $c$ and $c_1(\kappa)$ not depending on $n$ such that the event
%
\begin{equation}\label{eqn:corevent}
 \Big| \ts\frac{\hat \Sigma_{j,j'}}{\sigma_j\sigma_{j'}}-\rho_{j,j'}\Big| \leq \frac{c_1(\kappa)\log(n)}{\sqrt n}
\end{equation}
holds with probability at least $1-cn^{-\kappa}$. 
\end{lemma}
\paragraph{Remark}The event in the lemma has been formulated to hold with probability at least $1-cn^{-\kappa}$, rather than $1-\frac{c}{n}$, in order to accommodate a union bound for proving Lemma~\ref{lem:cor}.


\proof
 Consider the $\ell_2$-unit vectors $u=\Sigma^{1/2}e_j/\sigma_j$ and $v=\Sigma^{1/2}e_{j'}/\sigma_{j'}$  in $\R^p$.
Letting $W_n$ be as defined in~\eqref{eqn:Wndef}, observe that
\begin{equation}\label{eqn:decomp}
\begin{split}
\frac{\hat \Sigma_{j,j'}}{\sigma_j\sigma_{j'}} -\rho_{j,j'}
 & \ = u\ttop (W_n-\mathbf{I}_p)v.
%
%
\end{split}
\end{equation}
For each $1\leq i\leq n$, define the random variables $\zeta_{i,u}=Z_i\ttop u$ and $\zeta_{i,v}=Z_i\ttop v$. In this notation, the relation~\eqref{eqn:decomp} becomes
\begin{equation*}
\frac{\hat \Sigma_{j,j'}}{\sigma_j\sigma_{j'}} -\rho_{j,j'} =\underbrace{\Big(\ts\frac 1n \sum_{i=1}^n \zeta_{i,u}\zeta_{i,v}- u\ttop v\Big)}_{=:\Delta(u,v)} \ - \ \underbrace{\Big(\ts\frac 1n \sum_{i=1}^n \zeta_{i,u}\Big)\Big(\ts\frac 1n \sum_{i=1}^n \zeta_{i,v}\Big)}_{=:\Delta'(u,v)}.
\end{equation*}
Note that $\E[\zeta_{i,u}\zeta_{i,v}]=u\ttop v$.  Also, if we let $q=\max\{\kappa\log(n),3\}$, then
$$\|\zeta_{i,u}\zeta_{i,v}-u\ttop v\|_q \ \lesssim \ q^2,$$
which follows from Assumption~\ref{A:model}. Therefore, Lemma~\eqref{lem:rosenthal} gives the following bound for $q>2$,
\begin{equation}
\begin{split}
\|\Delta(u,v)\|_q & \ \lesssim  \ q \max\Big\{ \|\Delta(u,v)\|_2 \, , \, \ts\frac{1}{n}\big(\sum_{i=1}^n \|\zeta_{i,u}\zeta_{i,v}-u\ttop v\|_q^q\big)^{1/q}\Big\}\\[0.3cm]
%
& \ \lesssim \ q \max\Big\{\ts\frac{1}{\sqrt n} \, , \,  n^{-1+1/q}\cdot q^2\Big\}\\[0.3cm]
%
& \ \lesssim \ \ts\frac{\log(n)}{\sqrt n}.
\end{split}
\end{equation}
Using the Chebyshev inequality 
$$\P(|\Delta(u,v)|\geq e \|\Delta(u,v)\|_q) \ \leq e^{-q},$$
 we have
 $$\P\Big(|\Delta(u,v)|\geq \ts\frac{c\kappa\log(n)}{\sqrt n}\Big) \ \leq \ \ts\frac{1}{n^{\kappa}}.$$
Similar reasoning leads to the following tail bound for $\Delta'(u,v)$, 
  $$\P\bigg(|\Delta'(u,v)|\geq \big(\ts\frac{c\kappa\log(n)}{\sqrt n}\big)^2\bigg) \ \leq \ \ts\frac{1}{n^{\kappa}},$$
and combining with the previous tail bound gives the stated result.\qed
 
 
 


\begin{lemma}\label{lem:coupling}
Fix any number $\delta\in(0,1/2)$, and suppose the conditions of Theorem~\ref{THM:G} hold. Then, there is a constant $c>0$ not depending on $n$ such that the event~\eqref{eqn:maincouple} holds with probability at least $1-\frac cn$.
\end{lemma}

\proof Let $(a_1,\dots,a_{k_n})$ and $(b_1,\dots,b_{k_n})$ be real vectors, and note the basic fact
$$\bigg|\max_{1\leq j\leq k_n} a_j -\max_{1\leq j\leq k_n}b_j\bigg| \ \leq \ \max_{1\leq j\leq k_n} |a_j-b_j|.$$
From this, it is simple to derive the inequality
\begin{equation}
\big| \breve M_{k_n}^{\star}-M_{k_n}^{\star}\big| \leq \max_{j\in\J(k_n)}\Big|\big(\ts\frac{\hat\sigma_j}{\sigma_j}\big)^{\tau_n}-1\Big|\cdot \displaystyle \max_{j\in\J(k_n)}\Big|S_j^{\star}/\hat\sigma_j^{\tau_n}\Big|.
\end{equation}
To handle the first factor on the right side, it follows from Lemma~\ref{lem:corbasic} that the event
\begin{equation}
\max_{j\in\J(k_n)}\Big|\big(\ts\frac{\hat\sigma_j}{\sigma_j}\big)^{\tau_n}-1\Big| \ \leq c\cdot n^{-1/2}\cdot \log(n)
\end{equation}
holds with probability at least $1-\frac cn$. Next,  consider the random variable
\begin{equation}\label{eqn:Ustar}
U^{\star}:=\max_{j\in\J(k_n)} |S_{n,j}^{\star}/\hat\sigma_j^{\tau_n}|.
\end{equation}
It suffices to show there is possibly larger constant $c>0$, such that the event
\begin{equation}\label{eqn:Ustarbound}
\P\Big( U^{\star} \geq  c \log(n)^{3/2} \Big| X\Big) \leq \ts\frac 1n 
\end{equation}
holds with probability at least $1-\frac cn$.
%
Using Chebyshev's inequality  with $q\geq \log(n)$ gives
 $$\P\Big(U^{\star} \geq e\, \big(\E[|U^{\star}|^q|X]\big)^{1/q}\Big| X\Big)\leq e^{-q}.$$
Likewise, if the event
 \begin{equation}\label{eqn:UstarLq}
 (\E[|U^{\star}|^q|X])^{1/q} \ \leq c\log(n)^{3/2}
 \end{equation}
holds for some constant $c>0$, then the event~\eqref{eqn:Ustarbound} also holds. For this purpose, the argument in the proof of Lemma~\ref{lem:Bboundboot}\eqref{bbound} can be essentially repeated with $q=\max\{\frac{2}{\beta_n},\log(n),3\big\}$ to show that the event~\eqref{eqn:UstarLq} holds with probability at least $1-\frac cn$. The main detail to notice when repeating the argument is that $U^{\star}$ involves a maximum over $\J(k_n)$, whereas the argument for Lemma~\ref{lem:Bboundboot}\eqref{bbound} involves a maximum over $\J(k_n)^c$. This distinction can be handled by using the bound~\eqref{eqn:boundmathfrakt} in Lemma~\ref{lem:mathfrak}.\qed


\paragraph{The case when \lowercase{ $p\leq k_n$}}\label{lowdimcase}
The previous proofs relied on the condition $p>k_n$ only insofar as this implies $k_n\geq n^{\frac{1}{\log(n)^a}}$ and $\ell_n\geq \log(n)^3$. (These conditions are used in the analyses of $\I_n$ and $\III_n$, as well as $\I_n'$ and $\III_n'(X)$.)  However, if $p\leq k_n$, then the definition of $k_n$ implies that $p=k_n$, which causes the quantities $\I_n$, $\III_n$, $\I_n'(X)$ and $\III_n'(X)$ to become exactly 0. In this case, the proofs of Theorems~\ref{THM:G} and~\ref{THM:BOOT} reduce to bounding $\II_n$ and $\II_n'(X)$, and these arguments can be repeated as before.

%
%
%
%
%
%
%
%
%
%
%
%
%
%

\section{Proof of Theorem~\ref{THM:M}}\label{app:mult}
\paragraph{Notation and remarks}  
%
An important piece of notation for this appendix (and the next one) is the integer $d_n$, which we define to be the largest index $d_n\in\{1,\dots,p\}$ such that $\sigma_{(d_n)}^2\geq \e_0^2/\sqrt{n}$, with $\e_0$ as in Assumption~\ref{A:mult}. (Such an index must exist under Assumption~\ref{A:mult}.)
Also, it is simple to check that $k_n\leq d_n$ holds for all large $n$ under Assumption~\ref{A:mult}. Furthermore, as in the previous appendices, we will assume $p>k_n$, since the case $p\leq k_n$ can be handled using similar reasoning to that explained in the previous paragraph. Lastly, it will be helpful to note that $\sigma_{(j)}^2=\pi_{(j)}(1-\pi_{(j)})$, since it can be checked that $\pi_i\leq \pi_j$ implies $\sigma_i\leq \sigma_j$.




\paragraph{Outline of proof} Since the number $d_n$ will often play the role that $p$ did in previous proofs, we will use a slightly different notation for the analogues of the earlier quantities $\I_n$, $\II_n$, $\II_n'(X)$, and $\III_n'(X)$. This will also serve as a reminder that new details are involved in the context of the multinomial model. The new quantities are:
\begin{align*}
\bI_n &= d_{\text{K}}\big(\mathcal{L}(M_{d_n}) \, , \, \mathcal{L}(M_{k_n})\big)\\[0.2cm]
%
\bII_n &=d_{\text{K}}\big( \mathcal{L}(M_{k_n})\, , \, \mathcal{L}(\tilde M_{k_n})\big)\\[0.2cm]
%
%
\bII'_n(X) &\ = \ d_{\text{K}}\big(\mathcal{L}(\tilde M_{k_n}) \, , \, \mathcal{L}( M^{\star}_{k_n}|X\big)\big)\\[0.2cm]
%
\bIII'_n(X) & \ = \ d_{\text{K}}\big(\mathcal{L}(M^{\star}_{k_n}| X\big) \, , \, \mathcal{L}( M_{d_n}^{\star} |X\big)\big).
\end{align*}
The overall structure of the proof is based on the simple bound
\begin{equation*}\label{eqn:threepartsM}
\begin{split}
   d_{\textup{K}}\big(\L(\M)\,,\,\L(\M^{\star}|X)\big) \ \leq & \  \ d_{\textup{K}}\big(\L(\M) \, ,\,\L(M_{d_n})\big)\\[0.2cm]
   %
   & \ + \bI_n +\bII_n+\bII_n'(X)+\bIII_n'(X)\\[0.2cm]
   %
   & \ + d_{\textup{K}}\big(\L(M_{d_n}^{\star}|X) \, , \, \L(\M^{\star}|X)\big).
   \end{split}
\end{equation*}
Most of the proof will be completed through four separate lemmas, showing that each of the quantities in Roman numerals is at most of order $n^{-1/2+\delta}$. These lemmas are labeled as~\ref{lem:IIM} (for $\bII_n$),~\ref{lem:IM} (for $\bI_n$),~\ref{lem:IIprimeM} (for $\bII'_n(X)$) and~\ref{lem:IIIprimeM} (for $\bIII_n'(X)$). Finally, the other two quantities in the first and third lines %
will be shown to be at most of order $n^{-1/2+\delta}$ in Lemma~\ref{lem:otherM}. \qed
~\\


\begin{lemma}\label{lem:IIM}Fix any number $\delta\in(0,1/2)$, and suppose the conditions of Theorem~\ref{THM:M} hold. Then,
\begin{equation}
\bII_n \ \lesssim \  n^{-1/2+\delta}.
\end{equation}
\end{lemma}
\proof The argument is rougly similar to the proof of Proposition~\ref{prop:bentkus}, and we retain some of the notation used there. For the new proof, define the $k_n\times k_n$ matrix 
\begin{equation}\label{eqn:sigmakdef}
\Sigma(k_n):=\Pi_{k_n}\Sigma\Pi_{k_n}\ttop,
\end{equation}
where $\Pi_{k_n}\in\R^{k_n\times p}$ is the projection onto the coordinates indexed by $\J(k_n)$, as explained on page~\pageref{eqn:IInbound}. Let $r_n$ denote the rank of $\Sigma(k_n)$, and let 
$$\Sigma(k_n) = Q\Lambda_{r_n} Q\ttop$$
be a spectral decomposition of $\Sigma(k_n)$, where $Q\in\R^{k_n\times r_n}$ has orthonormal columns, and $\Lambda\in\R^{r_n\times r_n}$ is diagonal and invertible. In addition, define
$$\mathfrak{C}=\Lambda_{r_n}^{1/2}Q\ttop D_{k_n}^{-\tau_n},$$
as well as the $r_n$-dimensional random vector
$$Z_i'=\Lambda_{r_n}^{-1/2}Q\ttop \Pi_{k_n}(X_i-\boldsymbol\pi),$$
and the sample average
$$\breve Z_n'= \ts\frac{1}{\sqrt n}\sum_{i=1}^n Z_i'.$$
Since the random vector $\Pi_{k_n}(X_i-\boldsymbol \pi)$ lies in the column span of $\Sigma(k_n)$ almost surely, it follows that the relation $QQ\ttop \Pi_{k_n}(X_i-\boldsymbol\pi)=\Pi_{k_n}(X_i-\boldsymbol\pi)$ holds almost surely. In turn, this gives
$$Q\Lambda_{r_n}^{1/2} Z_i' = \Pi_{k_n}(X_i-\boldsymbol \pi),$$
and hence
%
 $$D_{k_n}^{-\tau_n}\Pi_{k_n}S_n = \mathfrak{C}\ttop\breve Z_n'.$$
 Consequently, for any $t\in\R$, there is a Borel convex set $\mathcal{A}_t\subset\R^{r_n}$ such that 
$$\P\big(M_{k_n}\leq t) \ = \  \P\Big(\max_{j\in\J(k_n)}D_{k_n}^{-\tau_n}\Pi_{k_n}S_{n,j}  \leq t\Big) \ = \  \P\big(\breve Z_n'\in\mathcal{A}_t\big).$$
Similar reasoning also can be applied to $\tilde M_{k_n}$ in order to obtain the expression \smash{$\P(\tilde M_{k_n}\leq t)= \gamma_{r_n}(\mathcal{A}_t)$}, where $\gamma_{r_n}$ is the standard Gaussian distribution on $\R^{r_n}$. 
%
Therefore, we have the bound
\begin{equation}
\bII_n \ \leq  \ \sup_{\mathcal{A}\in\mathscr{A}}\Big| \P(\breve Z' \in\mathcal{A}) - \gamma_{r_n}(\mathcal{A})\Big|,
\end{equation}
with $\mathscr{A}$ being the collection of Borel convex subsets of $\R^{r_n}$.

Since the vectors $Z_1',\dots,Z_n'$ are i.i.d., with mean 0 and identity covariance matrix, we may apply the Berry-Esseen bound of Bentkus~(Lemma~\ref{lem:bentkus}). The only remaining detail is to bound $\E[\|Z_1'\|_2^3]$, and show that it is at most a fixed power of $k_n$. To do this, Lemma~\ref{lem:lambdamin} implies there is a constant $c>0$ not depending on $n$ such that the following inequalities hold with probability 1,
\begin{equation*}
\begin{split}
\|Z_1'\|_2^2  & \ \leq \ \ts\frac{1}{\lambda_{r_n}(\Sigma(k_n))}\big\|\Pi_{k_n}(X_i-\boldsymbol\pi)\big\|_2^2\\[0.2cm]
%
& \ \leq \ ck_n^c \cdot k_n,
\end{split}
\end{equation*}
where we have used the facts that $\|Q\|\op\leq 1$, and the entries of $\Pi_{k_n}(X_i-\boldsymbol\pi)$ are bounded in magnitude by 1. Thus, we have $\E[\|Z_1'\|_2^3]\lesssim k_n^{(3/2)(c+1)}$, which completes the proof.\qed


\begin{lemma}\label{lem:IM}
Fix any number $\delta\in(0,1/2)$, and suppose the conditions of Theorem~\ref{THM:M} hold. Then,
\begin{equation}
\bI_n \ \lesssim \ n^{-1/2+\delta}.
\end{equation}
\end{lemma}
\proof By repeating the proof of Proposition~\ref{prop:IandIII}, it follows that if $\texttt{t}_{1,n}$ and $\texttt{t}_{2,n}$ are any two real numbers with $\texttt{t}_{1,n}\leq \texttt{t}_{2,n}$, then 
\begin{equation}
\bI_n \ \leq \ \P(\texttt{A}(\texttt{t}_{2,n})) \ + \ \P(\texttt{B}(\texttt{t}_{1,n})),
\end{equation}
where we write
\small
\begin{equation*}
\texttt{A}(t)=\Big\{\max_{j\in\J(k_n)} S_{n,j}/\sigma_j^{\tau_n} \leq t\Big\} \ \ \ \text{ and } \ \  \ \ \texttt{B}(t)=\Big\{\max_{j\in\J(d_n)\setminus\J(k_n)} S_{n,j}/\sigma_j^{\tau_n}> t\Big\},
\end{equation*}
\normalsize
for any $t\in \R$. (Note that $\texttt{B}(t)$ differs from $B(t)$ only insofar as $\J(d_n)\setminus\J(k_n)$ replaces $\J(k_n)^c$.)

To handle the probability $\P(\texttt{A}(\texttt{t}_{2,n}))$, we mimic the definition~\eqref{eqn:t2} and let
\begin{equation}
 \texttt{t}_{2,n}=\e_0\cdot \ell_n^{-\beta_n}\cdot \sqrt{\log(\ell_n)},  
\end{equation}
with $\e_0\in(0,1)$ as in Assumption~\ref{A:mult}. Having made this choice, the proof of Lemma~\ref{lem:bounds}\eqref{aboundlem} may be repeated essentially verbatim to show that $\P(\texttt{A}(\texttt{t}_{2,n}))\lesssim n^{-1/2+\delta}$. In particular, it is important to note that the correlation matrix $R(\ell_n)$ in the multinomial case satisfies the conditions needed for that argument to work, because $R^+(\ell_n)=\mathbf{I}_{\ell_n}$. Also, this argument relies on Lemma~\ref{lem:IIM} in the same way that the proof of Lemma~\ref{lem:bounds}\eqref{aboundlem} relies on Proposition~\ref{prop:bentkus}.

To handle $\P(\texttt{B}(\texttt{t}_{1,n}))$, the proof of Lemma~\ref{lem:bounds}\eqref{bboundlem} can be mostly repeated to show that this probability is of order $1/n$. However, there are a few differences. First, we may regard the $\alpha$ from the context of Lemma~\ref{lem:bounds}\eqref{bboundlem} as being equal to $1/2$, due to the basic fact that $\sigma_{(j)}\leq  j^{-1/2}$ always holds in the multinomial model. Likewise, we define
 \begin{equation}
 \texttt{t}_{1,n}=c\cdot k_n^{-(1-\tau_n)/2}\cdot \log(n)  
\end{equation}
as the analogue of $t_{1,n}$ in~\eqref{eqn:t1}. The only other issues to notice are that the set $\J(k_n)^c$ is replaced by $\J(d_n)\setminus\J(k_n)$, and that we must verify the following condition. Namely, there is a constant $c>0$ not depending on $n$ such that the inequality
 $$\max_{j\in\J(d_n)} \big\|\ts\frac{1}{\sigma_j}S_{n,j}\big\|_q \ \leq c q$$
%
 holds when $q=\max\{\frac{2}{(1/2)(1-\tau_n)},\log(n),3\}$. This will be verified later in Lemma~\ref{lem:SnjnormM}. Finally, it is simple to check that  $\texttt{t}_{1,n}  \leq \texttt{t}_{2,n}$ holds for all large $n$. \qed
~\\

\begin{lemma}\label{lem:IIprimeM}
Fix any number $\delta\in(0,1/2)$, and suppose the conditions of Theorem~\ref{THM:M} hold. Then, there is a constant $c>0$ not depending on $n$ such that the event
\begin{equation}
\bII_n'(X) \ \leq \ c\, n^{-1/2+\delta}
\end{equation}
holds with probability at least $1-\frac cn$.
\end{lemma}

\proof Let the random variable $\breve M_{k_n}^{\star}$ be as defined in the proof of Proposition~\ref{prop:IIprime}, and consider the triangle inequality
 \begin{equation}\label{eqn:IIprimeM}
 \small
\bII_n'(X) \ \leq \ \dK\Big(\L(\tilde M_{k_n})\, , \L(\breve M_{k_n}^{\star}|X)\Big) \ + \ \dK \Big(\L(\breve M_{k_n}^{\star}|X) \, , \, \L(M_{k_n}^{\star}|X)\Big).
\end{equation}
Regarding the second term on the right, the proof of Proposition~\ref{prop:IIprime} shows that this term can be controlled with a coupling inequality for $\breve{M}_{k_n}^{\star}$ and $M_{k_n}^{\star}$, as well as an anti-concentration inequality for $M_{k_n}^{\star}$. In the multinomial context, both of these inequalities can be established using the same overall approach as before. The main items that need to be updated are to replace $\J(k_n)^c$ with $\J(d_n)\setminus \J(k_n)$, and to control quantities involving \smash{$\{\hat\sigma_j | j\in\J(d_n)\}$} by using Lemmas~(\ref{lem:corM},~\ref{lem:sighatrnormM}) instead of Lemmas~(\ref{lem:cor},~\ref{lem:sighatrnorm}).

Next, to control the first term in the bound~\eqref{eqn:IIprimeM}, we will use an argument based on  Lemma~\ref{lem:hellinger}, which requires quite a few pieces of notation. First, let
$\Sigma(k_n)=\Pi_{k_n}\Sigma\Pi_{k_n}\ttop \in\R^{k_n\times k_n}$, and let the rank of this matrix be denoted by $r_n$. Next, write the spectral decomposition of $\Sigma(k_n)$ as
 $$\Sigma(k_n)=Q\Lambda_{r_n}Q\ttop,$$
 where $Q\in\R^{k_n\times r_n}$ has orthonormal columns and $\Lambda_{r_n}\in\R^{r_n\times r_n}$ is diagonal. In addition define 
\begin{equation}\label{eqn:whitemultdef}
\begin{split}
\hat\Sigma(k_n)&=\Pi_{k_n}\hat\Sigma \Pi_{k_n}\ttop\\[0.2cm]
W_n&= \Lambda_{r_n}^{-1/2}Q\ttop \hat\Sigma(k_n) Q\Lambda_{r_n}^{-1/2}.
\end{split}
\end{equation}
Since each vector $\Pi_{k_n}(X_i-\bar X)$ lies in the column span of $\Sigma(k_n)$ almost surely, it follows that the relation $QQ\ttop\hat\Sigma(k_n)QQ\ttop=\hat\Sigma(k_n)$ holds almost surely, which is equivalent to
$$\hat\Sigma(k_n)=Q\Lambda_{r_n}^{1/2}W_n\Lambda_{r_n}^{1/2}Q\ttop.$$
With this in mind,  let $\mathfrak{C}=\Lambda_{r_n}^{1/2}Q\ttop D_{k_n}^{-\tau_n}$, and also define
\begin{equation}
\begin{split}
\mathfrak{S}&=\mathfrak{C}\ttop \mathfrak{C}\\[0.2cm]
\breve{\mathfrak{S}} &= \mathfrak{C}\ttop W_n\mathfrak{C}.
\end{split}
\end{equation}
It is straightforward to check that $\tilde M_{k_n}$ is the coordinate-wise maximum of a Gaussian vector drawn from $N(0,\mathfrak{S})$, and similarly, $\breve M_{k_n}^{\star}$ is the coordinate-wise maximum of a Gaussian vector drawn from $N(0,\breve{\mathfrak{S}})$.

We will now compare $\tilde M_{k_n}$ and $\breve M_{k_n}^{\star}$ by applying Lemma~\ref{lem:hellinger}. For this purpose, let
\begin{equation}\label{eqn:Udef}
\mathfrak{C}=ULV\ttop
\end{equation}
be an s.v.d.~for $\mathfrak{C}$, where the matrix $L\in\R^{r_n\times r_n}$ is diagonal and invertible, and the matrices $U\in\R^{r_n\times r_n}$ and $V\in\R^{k_n\times r_n}$ each have orthonormal columns. From these definitions, it follows that
%
%
\begin{equation}
\Big(V\ttop \mathfrak{S} V\Big)^{-1/2}\Big( V\ttop \breve{\mathfrak{S}}V\Big) \Big(V\ttop\mathfrak{S}V\Big)^{-1/2} \ = \ U\ttop W_n U.
\end{equation}
Thus, the matrix $(V\ttop \mathfrak{S} V)^{1/2}$ will play the role of $H$ in the statement of Lemma~\ref{lem:hellinger}.
Also, in order to apply that lemma, we need that the columns of $\mathfrak{S}$ and $\breve{\mathfrak{S}}$ span the same subspace of $\R^{k_n}$ (with high probability). Noting that $\mathfrak{S}=VL^2 V\ttop$ and $\breve{\mathfrak{S}} =VL (U\ttop W_n U)L V\ttop$, it follows that $\mathfrak{S}$ and $\breve{\mathfrak{S}}$ have the same column span whenever $U\ttop W_n U$ is invertible, and by Lemma~\ref{lem:whitemult}, this holds with probability at least $1-c/n$. Therefore, Lemma~\ref{lem:hellinger} ensures that if the event
 \begin{equation}\label{eqn:tempevent}
 \|U\ttop W_n U-\mathbf{I}_{r_n}\|_{\text{op}} \ \leq \ \e,
 \end{equation}
 holds for some number $\e>0$, then the event
\begin{equation}
 \dK\Big(\L(\tilde M_{k_n})\, , \L(\breve M_{k_n}^{\star}|X)\Big) \leq c\cdot  k_n^{1/2} \cdot \e
\end{equation}
also holds. Finally, Lemma~\ref{lem:whitemult} shows that if we take $\e= c\log(n)k_n^c/\sqrt{n}$, then the event~\eqref{eqn:tempevent} holds with probability at least $1-c/n$, which completes the proof.\qed




\begin{lemma}\label{lem:IIIprimeM}
Fix any number $\delta\in(0,1/2)$, and suppose the conditions of Theorem~\ref{THM:M} hold. Then, there is a constant $c>0$ not depending on $n$ such that the event
\begin{equation}
\bIII_n'(X) \ \leq \ c\, n^{-1/2+\delta}
\end{equation}
holds with probability at least $1-\frac cn$.
\end{lemma}

\proof The proof follows the argument outlined in Section~\ref{sec:III'}. The main details to be updated for the multinomial context arise in controlling the quantities $\{\hat\sigma_j | j\in\J(d_n)\}$. Specifically, the index set $\J(k_n)^c$ must be replaced with  $\J(d_n)\setminus \J(k_n)$, and Lemmas~(\ref{lem:corM},~\ref{lem:sighatrnormM}) must be used in place of Lemmas~(\ref{lem:cor},~\ref{lem:sighatrnorm}).\qed



\begin{lemma}\label{lem:otherM}
Suppose the conditions of Theorem~\ref{THM:M} hold. Then, there is a constant $c>0$ not depending on $n$, such that
\begin{equation}\label{eqn:reductionfirst}
 d_{\textup{K}}\big(\L(\M), \L(M_{d_n})\big) \ \leq \ \bI_n+\ts\frac{c}{n},
\end{equation}
and the event
\begin{equation}\label{eqn:reductionsecond}
  d_{\textup{K}}\big(\L(\M^{\star}|X), \L(M_{d_n}^{\star}|X)\big) \ \leq \   \bIII_n'(X) 
\end{equation}
occurs with probability at least $1-\frac cn$.
\end{lemma}
%
%
%
%
%
%
%
%
%
\proof Fix any $t\in\R$. By intersecting the event $\{\max_{j\in\hat\J_n}S_{n,j}/\sigma_j^{\tau_n}\leq t\}$ with the events $\{\J(k_n)\subset \hat\J_n\}$ and $\{\J(k_n)\not \subset\hat\J_n\}$, and noting that the maximum can only become smaller on a subset, we have
\begin{equation*}
\footnotesize
\begin{split}
\P\Big(\max_{j\in\hat\J_n}S_{n,j}/\sigma_j^{\tau_n}\leq t\Big) \ \leq \
%
 \P\Big(\max_{j\in\J(k_n)}S_{n,j}/\sigma_j^{\tau_n}\leq t\Big) +\P\big(\J_{k_n}\not\subset\hat\J_n).
\end{split}
\end{equation*}
Therefore, by subtracting from both sides the probability involving the maximum over $\J(d_n)$, we have
\footnotesize
\begin{equation}\label{eqn:firstsign}
 \P\Big(\max_{j\in\hat\J_n}S_{n,j}/\sigma_j^{\tau_n}\leq t\Big) - \P\Big(\max_{j\in\J(d_n)}S_{n,j}/\sigma_j^{\tau_n}\leq t\Big) \ \leq \ \bI_n+\P\big(\J_{k_n}\not\subset\hat\J_n).
\end{equation}
\normalsize
%
Similarly, by intersecting with the events $\{\hat\J\subset \J(d_n)\}$ and $\{\hat\J_n\not \subset\J(d_n)\}$, we have
\footnotesize
\begin{equation*}
%\begin{split}
\P\Big(\max_{j\in\J(d_n)}S_{n,j}/\sigma_j^{\tau_n}\leq t\Big) 
 \ \leq \  \P\Big(\max_{j\in\hat\J_n}S_{n,j}/\sigma_j^{\tau_n}\leq t\Big)+ \P(\hat\J_n\not \subset\J(d_n)).
%\end{split}
\end{equation*}
\normalsize
 Next, subtracting the probability involving the maximum over $\hat\J_n$ gives
\footnotesize
\begin{equation}\label{eqn:secondsign}
%\begin{split}
\P\Big(\max_{j\in\J(d_n)}S_{n,j}/\sigma_j^{\tau_n}\leq t\Big) 
 \   - \ \P\Big(\max_{j\in\hat\J_n}S_{n,j}/\sigma_j^{\tau_n}\leq t\Big) \ \leq \ 0 + \P(\hat\J_n\not \subset\J(d_n)).
%\end{split}
\end{equation}
\normalsize
Combining~\eqref{eqn:firstsign} and~\eqref{eqn:secondsign} implies
\begin{equation*}
d_{\textup{K}}\big(\L(\M)\, ,\, \L(M_{d_n})\big) \ \leq  \ \I_n+\P\big(\J_{k_n}\not\subset\hat\J_n)+\P(\hat\J_n\not \subset\J(d_n)),
\end{equation*}
and furthermore, Lemma~\ref{lem:Jhat} shows that the last two probabilities on the right are at most $c/n$. This proves ~\eqref{eqn:reductionfirst}. The inequality~\eqref{eqn:reductionsecond} follows by similar reasoning, and is actually easier, because conditioning on $X$ allows us to work under the assumption that the events $\{\J(k_n)\subset \hat\J_n\}$ and $\{\hat\J_n\subset\J(d_n)\}$ hold, since they occur with probability at least $1-\frac cn$. \qed

%
%
%
%
%
%
%
%
%
%
%
%
%

\section{Technical Lemmas for Theorem~\ref{THM:M}}\label{app:lemmamult}


\begin{lemma}\label{lem:Jhat}
Suppose the conditions of Theorem~\ref{THM:M} hold. Then, with probability at least $1-\frac cn$, the following two events hold simultaneously,
\begin{equation}\label{eqn:firstJhat}
\J(k_n) \ \subset \ \hat{\mathcal{J}}_n, 
\end{equation}
and
\begin{equation}\label{eqn:secondJhat}
\hat{\mathcal{J}}_n \ \subset \ \J(d_n).
\end{equation}
\end{lemma}
%%%%%%%%%%%%%%
\proof We first address the event~\eqref{eqn:firstJhat}. By a union bound, the following inequalities hold for all large $n$,
\begin{equation}
\begin{split}
\P\Big(\J(k_n)\not\subset \hat{\mathcal{J}}_n \Big) 
%
& \ \leq \  \sum_{j\in\J(k_n)} \P\Big(\hat\pi_j <\sqrt{\log(n)/n}\Big)\\[0.3cm]
%
%
&  \leq \  k_n\cdot \max_{j\in\J(k_n)}\P\Big(\sqrt{n}|\hat\pi_j-\pi_j| >n^{1/4}\Big)
\end{split}
\end{equation}
where the last step follows from the fact that if $j\in\J(k_n)$, then the crude (but adequate) inequality $\sqrt{n}\pi_j- \sqrt{\log(n)}\geq n^{1/4}$ holds for all large $n$. In turn, Hoeffding's inequality~\cite[p.460]{vanderVaart:Wellner:2000} implies
%
$$ k_n\cdot\max_{j\in\J(k_n)}\P\Big(\sqrt{n}|\hat\pi_j-\pi_j| >n^{1/4}\Big) \ \leq \ 2k_n e^{-2\sqrt{n}} \ \lesssim \ \ts\frac{1}{n},$$
which establishes~\eqref{eqn:firstJhat}.

We now turn to~\eqref{eqn:secondJhat}. Under Assumption~\ref{A:mult}, it is simple to check that $\J(k_n)\subset\J(d_n)$ holds for all large $n$. Consequently, in the low-dimensional situation where $k_n=p$, we must have $\J(d_n)=\J(p)$ for all large $n$, and then \eqref{eqn:secondJhat} is clearly true. Hence, we may work in the situation where $k_n<p$. To  begin, observe that Assumption~\ref{A:mult} gives $\pi_{(1)}(1-\pi_{(1)})=\sigma_{(1)}^2\geq \e_0^2$, which implies $1-\pi_{(1)}\geq \e_0^2$, and hence $1-\pi_j\geq \e_0^2$ for all $j\in\{1,\dots,p\}$. Based on this observation, if we consider the set 
$$\J_n':=\Big\{j\in\{1,\dots,p\}\, \Big| \, \pi_j\geq 1/\sqrt{n}\Big\},$$
then $j\in\J_n'$ implies $\sigma_j^2=\pi_j(1-\pi_j)\geq \e_0^2/\sqrt{n}$. Now, recall that $\J(d_n)$ is defined so that $j\in\J(d_n) \Longleftrightarrow \sigma_j^2\geq \e_0^2/\sqrt{n}$. As a result of this definition, we have $\J_n'\subset\J(d_n)$. Therefore, in order to show that the event  $\{\hat\J_n\subset\J(d_n)\}$ holds with probability at least $1-c/n$, it suffices to show that the event $\{\hat\J_n\subset\J_n'\}$ holds with at least the same probability.

To proceed,
observe that the following inclusion always holds,
\begin{equation}
\begin{split}
\hat{\mathcal{J}}_n
%
& \ \subset \ \Big\{ j\in\{1,\dots,p\} \ \Big| \  \sqrt{n}\pi_j \ \geq \sqrt{\log(n)} - \max_{1\leq j\leq p}\sqrt{n}(\hat\pi_j-\pi_j)\Big\}.
\end{split}
\end{equation}
Therefore, if we can show that the event 
$$\mathcal{E}:=\Big\{ \sqrt{\log(n)} - \max_{1\leq j\leq p}\sqrt{n}(\hat\pi_j-\pi_j) \geq 1\Big\}$$ 
occurs with probability at least $1-\frac cn$, then the event $\{\hat\J_n\subset \J_n'\}$ will also occur with probability at least $1-\frac cn$. Now, consider the union bound
\begin{equation}
\P(\mathcal{E}^c) \ \leq \  \sum_{j=1}^p \P\Big(\sqrt{n}(\hat\pi_j-\pi_j) > \sqrt{\log(n)}-1\Big).
\end{equation}
We will bound this sum by considering two different sets of indices. For the indices $j\in\J(k_n)^c$, the values $\pi_j$ are mostly small. This motivates the use of Kiefer's inequality (Lemma~\ref{lem:Kiefer}), which implies there is some $c>0$ not depending on $n$, such that the following bound holds for all large $n$, and $j\in\J(k_n)^c$,
\begin{equation}
\begin{split}
\P\Big(\sqrt{n}(\hat\pi_j-\pi_j) > \sqrt{\log(n)}-1\Big) & \ \leq \ 2\exp\Big\{-c\log(n)\log(\ts\frac{1}{\pi_j})\Big\}\\[0.2cm]
%
& \ = \ 2 \pi_j^{c\log(n)}.
%\\[0.3cm]
%
\end{split}
\end{equation}
On the other hand, if $j\in\J(k_n)$, then $\pi_j$ is of moderate size, and Hoeffding's inequality~\cite[p.460]{vanderVaart:Wellner:2000} implies the following inequality for all large $n$,
\begin{equation}
\begin{split}
\P\Big(\sqrt{n}(\hat\pi_j-\pi_j) > \sqrt{\log(n)}-1\Big) & \ \leq \ 2\exp\big\{-(3/2)\log(n)\big\}\\[0.2cm]
%
& \ = \ 2 n^{-3/2}.
%
\end{split}
\end{equation}
(The constant $3/2$ in the exponent has been chosen for simplicity, and is not of special importance.) Combining the two different types of bounds, and using the fact that any probability vector satisfies $\pi_{(j)}\leq j^{-1}$ for all $j\in\{1,\dots,p\}$, we obtain
\begin{equation}
\begin{split}
\P(\mathcal{E}^c) 
%
& \ \lesssim \ k_n n^{-3/2} \ + \ \sum_{j=k_n+1}^p \pi_{(j)}^{c\log(n)}\\[0.3cm]
%
& \ \lesssim \ \ts\frac{1}{n} \ +  \displaystyle\int_{k_n}^p x^{-c\log(n)}dx\\[0.3cm]
%
& \ \lesssim \ \ts\frac{1}{n} \ + k_n^{-c\log(n)+1}\\[0.3cm]
%
& \ \lesssim \ \ts\frac{1}{n},
\end{split}
\end{equation}
as needed.\qed




















\begin{lemma}\label{lem:sighatrnormM}
Suppose the conditions of Theorem~\ref{THM:M} hold, and let 
$q=\max\{\ts\frac{2}{(1/2)(1-\tau_n)}, \log(n),3\}$.
%
Then, there is a constant $c>0$ not depending on $n$, such that for any $j\in\mathcal{J}(d_n)$, we have
\begin{equation}
\|\hat\sigma_{j}\|_q \ \leq \  c \cdot \sigma_j \cdot \sqrt{q}.
\end{equation}
\end{lemma}
\proof By direct calculation
\begin{equation*}
%\label{eqn:sigsteps}
\begin{split}
\ts\frac{1}{\sigma_j}\|\hat\sigma_{j}\|_q
%
&= \ \ts\frac{1}{\sigma_j} \Bigg\|\ts\frac 1n \sum_{i=1}^n (X_{i,j}^2-\pi_j)+\big(\pi_j-\pi_j^2\big)+\big(\pi_j^2-\bar{X}_j^2\big)\Bigg\|_{q/2}^{1/2}\\[0.2cm]
%
&\leq \ \ts\frac{1}{\sigma_j} \Bigg(\big\|\ts\frac 1n \sum_{i=1}^n (X_{i,j}^2-\pi_j)\big\|_{q/2}^{1/2}+(\pi_j-\pi_j^2)^{1/2}+\big\|\pi_j^2-\bar{X}_j^2\big\|_{q/2}^{1/2}\Bigg)\\[0.3cm]
%
&\leq \ \ts\frac{1}{\sigma_jn^{1/4}}\big\| S_{n,j}\big\|_{q/2}^{1/2} \ + \ 1 \ + \ \ts\frac{\sqrt{2}}{\sigma_jn^{1/4}}\big\| S_{n,j}\big\|_{q/2}^{1/2}.
%
\end{split}
\end{equation*}
Since $1/(\sigma_j n^{1/4})\leq 1/\e_0$ when $j\in\J(d_n)$, it follows from Lemma~\ref{lem:rosenthal} that the first and third terms are at most of order $\sqrt{q}$. \qed


\paragraph{Remark} For the next lemma, put $\boldsymbol \pi_{k_n}:=(\pi_{(1)},\dots,\pi_{(k_n)})$, and recall the definition $\Sigma(k_n)=\Pi_{k_n}\Sigma\Pi_{k_n}\ttop$ from line~\eqref{eqn:sigmakdef}.


\begin{lemma}\label{lem:lambdamin}
%
If $r_n$ denotes the rank of $\Sigma(k_n)$, and the conditions of Theorem~\ref{THM:M} hold, then there is a constant $c>0$ not depending on $n$ such that
\begin{equation}
\lambda_{r_n}(\Sigma(k_n)) \ \gtrsim \ k_n^{-c}.
\end{equation}
%
\end{lemma}
\proof  We first consider the case $k_n<p$, and then handle the case $k_n=p$ separately at the end of the proof. 
 Under the multinomial model, we have for all $i,j\in\{1,\dots,k_n\}$,
$$\Sigma_{i,j}(k_n) = \begin{cases} \  \pi_{(i)}(1-\pi_{(i)}) & \text{ if } i=j, \\ \ -\pi_{(i)}\pi_{(j)} & \text{ if } i\neq j. \end{cases}$$
For each $i\in\{1,\dots,k_n\}$, define the ``deleted row sum'', 
$$\varrho_i:=\sum_{j\neq i}|\Sigma_{i,j}(k_n)|.$$
By the Ger\v{s}gorin disc theorem~\citep[Sec. 6.1]{Horn:Johnson:1990}, 
%
\begin{equation}
\begin{split}
\lambda_{r_n}(\Sigma(k_n)) & \ \geq \ \lambda_{\min}(\Sigma(k_n))\\[0.2cm]
 & \ \geq \  \min_{1\leq i\leq k_n}\Big\{\Sigma_{i,i}(k_n)  - \varrho_i\Big\}\\[0.2cm]
%
& \ = \ \min_{1\leq i\leq k_n} \Big\{(\pi_{(i)} -\pi_{(i)}\tsum_{j=1}^{k_n}\pi_{(j)}\Big\},\\[0.2cm]
%
& \ = \ \pi_{(k_n)}\Big(1-\tsum_{j=1}^{k_n} \pi_{(j)}\Big).
\end{split}
\end{equation}
When $k_n<p$, it follows from Assumption~\ref{A:mult} that
\begin{equation}
\begin{split}
1-\tsum_{j=1}^{k_n}\pi_{(j)}&\geq \pi_{(k_n+1)}\\[0.2cm]
%
&\geq \sigma_{(k+1)}^2\\[0.2cm]
%
&\geq \e_0^2 (k_n+1)^{-2\alpha}\\[0.2cm]
%
&\gtrsim k_n^{-2\alpha}.
\end{split}
\end{equation}
Hence, the previous steps show that $\lambda_{r_n}(\Sigma(k_n))$ is at least of order $k_n^{-4\alpha}$.

Finally, consider the case when $k_n=p$. In this case, it is a basic fact that the rank of $\Sigma=\Sigma(k_n)$ satisfies $r_n=p-1$ (where we note that  Assumption~\ref{A:mult} ensures $\pi_{(p)}>0$). Also, it is known from matrix analysis~\citep[Theorem 1]{Benasseni:2012} that 
$$\lambda_{p-1}(\Sigma)\geq \pi_{(p)},$$
and therefore Assumption~\ref{A:mult} leads to
\begin{equation}
\begin{split}
\lambda_{r_n}(\Sigma(k_n))&\geq \pi_{(k_n)}\\[0.2cm]
%
&\geq\sigma_{(k_n)}^2\\[0.2cm]
%
&\geq \e_0^2 k_n^{-2\alpha},
%
\end{split}
\end{equation}
as needed.\qed


\begin{lemma}\label{lem:whitemult} 
Let the deterministic matrix $U\in\R^{r_n\times r_n}$ and the random matrix $W_n\in\R^{r_n\times r_n}$ be as defined in~\eqref{eqn:Udef} and~\eqref{eqn:whitemultdef} respectively. Also, suppose that the conditions of Theorem~\ref{THM:M} hold. Then, there is a constant $c>0$ not depending on $n$ such that the event
\begin{equation}
\big\|U\ttop W_n U-\mathbf{I}_{r_n}\big\|_{\textup{op}}\ \leq \ \ts\frac{c\, k_n^c \log(n)}{\sqrt n}
\end{equation}
holds with probability at least $1-\frac cn$.
\end{lemma}
\proof Let the notation from the proof of Lemma~\ref{lem:IIprimeM} be in force, and observe that
\footnotesize
\begin{equation*}
\begin{split}
\big\|U\ttop W_n U-\mathbf{I}_{r_n}\big\|_{\textup{op}} & \ = \ \Big\|U\ttop\Big(\Lambda_{r_n}^{-1/2}Q\ttop \hat\Sigma(k_n) Q\Lambda_{r_n}^{-1/2} - \mathbf{I}_{r_n}\Big)U\Big\|\op\\[0.2cm]
%
& \ \leq \ \Big\|\Lambda_{r_n}^{-1/2}Q\ttop \Big(\hat\Sigma(k_n) -\Sigma(k_n)\Big)Q\Lambda_{r_n}^{-1/2}\Big\|\op\\[0.2cm]
%
& \ \leq \ \frac{1}{\lambda_{r_n}(\Sigma(k_n))} \, \big\|\hat\Sigma(k_n)-\Sigma(k_n)\big\|\op.
\end{split}
\end{equation*}
\normalsize
With regard to the first factor in the previous line, Lemma~\ref{lem:lambdamin} implies
$$\ts\frac{1}{\lambda_{r_n}(\Sigma(k_n))} \lesssim k_n^c.$$
To complete the proof, let $\boldsymbol\pi_{k_n}=\Pi_{k_n}\boldsymbol \pi$, and let $u\in\R^{k_n}$ be a generic unit vector. Consider the decomposition
\begin{equation*}
\footnotesize
\begin{split}
|u\ttop\big(\hat\Sigma(k_n)-\Sigma(k_n)\big)u| & \leq  \bigg|\frac1n\sum_{i=1}^n \Big((X_i\ttop \Pi_{k_n}\ttop u)^2-u\ttop\text{diag}(\boldsymbol\pi_{k_n})u\Big)\bigg| + \Big| (\bar X\ttop \Pi_{k_n}\ttop u)^2-(\boldsymbol\pi_{k_n}\ttop u)^2\Big|\\[0.3cm]
%
& \ =: \ \Delta_n(u)+\Delta_n'(u).
\end{split}
\end{equation*}
\normalsize
In order to control these terms, note that each random variable $(X_i\ttop \Pi_{k_n}\ttop u)^2$ is bounded in magnitude by 1, and has expectation equal to $u\ttop \text{diag}(\boldsymbol \pi_{k_n})u$. In addition, we have
$$|\Delta_n'(u)| \ \leq \ 2|\bar X\ttop \Pi_{k_n}\ttop u-\boldsymbol \pi_{k_n}\ttop u|.$$
Based on these observations, the proof of Lemma~\ref{lem:white} can be essentially repeated to show that there is a constant $c>0$ not depending on $n$ such that the event
$$ \big\|\hat\Sigma(k_n)-\Sigma(k_n)\big\|_{\text{op}} \ \leq \ \ts\frac{c\log(n) k_n}{\sqrt{n}}$$
holds with probability at least $1-\frac cn$. This completes the proof.\qed\\


\begin{lemma}\label{lem:SnjnormM} Let $q=\max\{\ts\frac{2}{(1/2)(1-\tau_n)}, \log(n),3\}$, and suppose that the conditions of Theorem~\ref{THM:M} hold.  Then, there is a constant $c>0$ not depending on $n$, such that 
\begin{equation}
\max_{j\in\mathcal{J}(d_n)}\|\ts\frac{1}{\sigma_j}S_{n,j}\|_q \leq cq.
\end{equation}
In addition, the following event holds with probability 1,
\begin{equation}
\max_{j\in\mathcal{J}(d_n)} \Big(\E\big[|\ts\frac{1}{\hat \sigma_j}S_{n,j}^{\star}|^q|X\big] \Big)^{1/q} \ \leq c\, q.
\end{equation}
\end{lemma}
\proof The second inequality can be obtained by repeating the proof of Lemma~\ref{lem:Snjnorm}, since $S_n^{\star}$ is still Gaussian under the setup of Assumption~\ref{A:mult}. To prove the first inequality, note that since $q>2$, Lemma~\ref{lem:rosenthal} gives
\begin{equation}\label{eqn:lemrosenthalfirst}
 \|\ts\frac{1}{\sigma_j}S_{n,j}\|_q \ \lesssim \ q\cdot \max\Big\{  \|\ts\frac{1}{\sigma_j}S_{n,j}\|_2 \, , \,  n^{-1/2+1/q} \|\ts\frac{1}{\sigma_j}(X_{1,j}-\pi_j)\|_q\Big\}.
\end{equation}
Clearly,
\begin{equation*}
 \|\ts\frac{1}{\sigma_j}S_{n,j}\|_2^2 \ = \  \var(\ts\frac{1}{\sigma_j}S_{n,j}) \ = \ 1.
\end{equation*}
For the stated choice of $q$, the second term inside the maximum satisfies
$$n^{-1/2+1/q} \|\ts\frac{1}{\sigma_j}(X_{1,j}-\pi_j)\|_q \ \lesssim \frac{1}{\sqrt n \sigma_j},$$
and also, since $j\in\mathcal{J}(d_n)$, we have $\frac{1}{\sqrt{n}\sigma_j}\lesssim \frac{1}{n^{1/4}}$, which leads to the stated claim. \qed



\begin{lemma}\label{lem:corM}
Suppose the conditions of Theorem~\ref{THM:M} hold. Then, there is a constant $c>0$ not depending on $n$ such that the events
\begin{equation}\label{eqn:firstcor}
 \max_{j\in\J(k_n)}\Big| \ts\frac{\hat\sigma_j}{\sigma_j}-1\Big| \leq \ts \frac{c\cdot k_n^c\cdot\sqrt{\log(n)}}{n^{1/2}},
\end{equation}
%
and
%
\begin{equation}\label{eqn:minvarhattemp}
 \min_{j\in\J(k_n)}\hat\sigma_j^{1-\tau_n} \ \geq \ \Big(\min_{j\in\J(k_n)} \sigma_j^{1-\tau_n}\Big)\cdot \Big(1-\ts\frac{c\cdot k_n^c\cdot \sqrt{\log(n)}}{n^{1/2}}\Big)
\end{equation}
each
hold with probability at least $1-\ts\frac{c}{n}$.
\end{lemma}

\proof  Note that if~\eqref{eqn:firstcor} occurs, then~\eqref{eqn:minvarhattemp} also occurs, and so we only deal with the former event. Fix any number $\kappa\geq 2$. By a union bound, it suffices to show there are positive constants $c$ and $c_1(\kappa)$ not depending on $n$, such that for any $j\in\J(k_n)$, the event
%
%
\begin{equation}\label{eqn:corevent}
 \Big| \ts\frac{\hat{\sigma}_j^2}{\sigma_j^2}-1\Big| \ \leq \ \frac{c_1(\kappa)\cdot k_n^c\cdot \sqrt{\log(n)}}{n^{1/2}}
\end{equation}
holds with probability at least $1-\frac{c}{n^{\kappa}}$. 
%
To this end, observe that
\begin{equation*}
\begin{split}
\Big| \ts\frac{\hat{\sigma}_j^2}{\sigma_j^2}-1\Big| & \ \leq \ \Big|\ts\frac{1}{\sigma_j^2n} \sum_{i=1}^n (X_{i,j}^2-\pi_j) \Big| \ + \ \ts\frac{1}{\sigma_j^2}\big|\pi_j^2-\bar{X}_j^2\big|\\[0.3cm]
 %
  & \ \leq \  \ts\frac{3}{\sigma_j^2 \sqrt{n}}|S_{n,j}|.
%
 \end{split}
\end{equation*}
Due to Assumption~\ref{A:mult}, and the fact that $j\in\J(k_n)$, we have 
$$\ts\frac{1}{\sigma_j^2\sqrt n} \ \lesssim  \ \ts\frac{k_n^{2\alpha}}{\sqrt n}.$$
In addition, Hoeffding's inequality ensures there is a constant $c_1(\kappa)$ such that the event
\begin{equation*}
|S_{n,j}| \ \leq \ c_1(\kappa)\sqrt{\log(n)}
\end{equation*}
holds with probability at least $1-\frac{c}{n^{\kappa}}$.\qed











%%%%%%%%%%%%%%%%%%%%%%%%%%%%%%%%%




\section{Background results} \label{app:background}
The following result is a multivariate version of the Berry-Esseen theorem due to~\cite{Bentkus:2003}. 
\begin{lemma}[Bentkus' multivariate Berry-Esseen theorem]\label{lem:bentkus}
\hspace{-0.4cm} Let\ $V_1,\dots,V_n$ be i.i.d.~random vectors $\R^d$, with zero mean, and identity covariance matrix. Furthermore, let $\gamma_d$ denote the standard Gaussian distribution on $\R^d$,  and let $\mathscr{A}$ denote the collection of all Borel convex subsets of $\R^d$. Then, there is an absolute constant $c>0$ such that
\begin{equation}
 \sup_{\mathcal{A}\in\mathscr{A}} \Big| \P\big(\ts\frac{1}{\sqrt{n}}(V_1+\dots+V_n)\in \mathcal{A}\big) -\gamma_d(\mathcal{A})\Big| \ \leq  \ \displaystyle \frac{c\cdot  d^{1/4}\cdot  \E\big[\|V_1\|_2^3\big]}{n^{1/2}}.
\end{equation}
\end{lemma}

\noindent The following is a version of Nazarov's inequality~\citep{Nazarov:2003,Klivans:2008}, as formulated in~\cite[Lemma 4.3]{CCK:SPA}.
\begin{lemma}[Nazarov's inequality]\label{lem:Nazarov}
Let $(\xi_1,\dots,\xi_m)$ be a multivariate normal random vector, and suppose the parameter $\underline{\sigma}^2:=\min_{1\leq j\leq m} \var(\xi_j)$ is positive. Then, for any $r>0$,
\begin{equation}
\sup_{t\in\R}\, \P\Big(\Big|\max_{1\leq j\leq m} \xi_j -t\Big| \ \leq r \Big)  \ \leq \  \frac{2r}{\underline{\sigma}}\cdot (\sqrt{2\log(m)}+2).
\end{equation}
\end{lemma}

\noindent The result below is a version of Slepian's lemma, which is adapted from~\cite[Theorem 2.2]{Li:Shao:2002}. (See references therein for earlier versions of this result.)
\begin{lemma}[Slepian's lemma]\label{lem:slepian} 
Let $m\geq 3$, and let $\mathsf{R}\in\R^{m\times m}$ be a correlation matrix with $\max_{i\neq j}\mathsf{R}_{i,j}<1$. Also, let $\mathsf{R}^+$ be the matrix with $(i,j)$ entry given by $\max\{\mathsf{R}_{i,j},0\}$, and suppose $\mathsf{R}^+$ is positive semi-definite. Furthermore, let $\zeta\sim N(0,\mathsf{R})$ and $\xi\sim N(0,\mathsf{R}^+)$. Then, the following inequalities hold for any $t\geq 0$,
\begin{equation}
 \P\Big(\max_{1\leq j\leq m} \zeta_j  \leq t\Big) \ \leq \ \P\Big(\max_{1\leq j\leq m} \xi_j\leq t\Big) \ \leq \  K_m(t) \cdot \Phi^m(t),
\end{equation}
where
\begin{equation}
K_m(t)=\exp\Bigg\{\sum_{1\leq i<j\leq m} \log\Big(\ts\frac{1}{1-\ts\frac{2}{\pi}\arcsin(\mathsf{R}_{i,j}^+)}\Big)\exp\Big(-\ts\frac{t^2}{1+\mathsf{R}_{i,j}^+}\Big)\Bigg\}.
\end{equation}

\end{lemma}

%
%
%
%
\noindent The following inequalities are due to~\cite{Zinn:1985}. 
%
%
\begin{lemma}[Rosenthal's inequality with best constants]\label{lem:rosenthal}
Fix $r\geq 1$ and put $\textup{Log}(r):=\max\{\log(r),1\}$. Let $\xi_1,\dots,\xi_m$ be independent random variables satisfying $\E[|\xi_j|^r]<\infty$ for all $1\leq j\leq m$. Then, there is an absolute constant $c>0$ such that the following two statements are true.
\begin{itemize}
\item[(i).] If $\xi_1,\dots,\xi_m$ are non-negative random variables, then 
	\begin{equation}
	\big\|\tsum_{j=1}^m \xi_i\big\|_r\leq c\cdot \ts\frac{r}{\textup{Log}(r)} \cdot \max\bigg\{ \big\|\tsum_{j=1}^m \xi_j\big\|_1 \, , \, \big(\tsum_{j=1}^m \big\|\xi_i\|_r^r\big)^{1/r}\bigg\}.
	\end{equation}
%
%
\item[(ii).] If $r>2$, and the random variables $\xi_1,\dots,\xi_m$ all have mean 0, then
	\begin{equation}
	\big\|\tsum_{j=1}^m \xi_i\big\|_r\leq c\cdot \ts\frac{r}{\textup{Log}(r)} \cdot \max\bigg\{ \big\|\tsum_{j=1}^m \xi_j\big\|_2 \, , \, \big(\tsum_{j=1}^m \big\|\xi_i\|_r^r\big)^{1/r}\bigg\}.
	\end{equation}
	\end{itemize}
	%
\end{lemma}

\paragraph{Remark} The non-negative case is handled in~\cite[Theorem 2.5]{Zinn:1985}.  With regard to the mean 0 case, the statement above differs slightly from \cite[Theorem 4.1]{Zinn:1985}, which requires symmetric random variables, but the 
 remark on page 247 of that paper explains why the variables $\xi_1,\dots,\xi_m$ need not be symmetric as long as they have mean 0.\\

The result below is a sharpened version of Hoeffding's inequality for handling the binomial distribution when the success probability is small~\citep[Corollary A.6.3]{vanderVaart:Wellner:2000}.
\begin{lemma}[Kiefer's inequality]\label{lem:Kiefer}
Let $\xi_1,\dots,\xi_m$ be independent Bernoulli random variables with success probability $\pi_0\in(0,1/e)$, and let $\bar\xi=\frac{1}{m}\sum_{i=1}^m\xi_i$. Then, the following inequality holds for any $m\geq 1$ and $t>0$,
\begin{equation}
\P\Big(\sqrt{m}|\bar\xi-\pi_0|\geq t\Big) \ \leq \ 2\exp\Big\{-t^2\big[\log\big(\ts\frac{1}{\pi_0}\big)-1\big]\Big\}.
\end{equation}
\end{lemma}
%%%%%%%%%%%%%%%%%%%%%%%%%%%%%%%%%




\section{Related work on Gaussian approximation}\label{app:Gaussian}

Although our focus is primarily on rates of bootstrap approximation, this topic is closely related to rates of Gaussian approximation in the central limit theorem --- for which there is a long line of work in finite-dimensional and infinite-dimensional settings. We refer to the chapter \citep{Bentkus:2000} for a general survey, as well as Appendix L of the paper~\citep{CCK:2013} for a discussion that is oriented more towards high-dimensional statistics.

 To describe how our work fits into this literature, we fix some notation. Let $\mathbb{B}$ denote a Banach space, with $\mathscr{A}$ being a collection of its subsets, and $\varphi:\mathbb{B}\to\R$ being a function. For the moment, we will regard the summands of $S_n$ as centered~i.i.d.~random elements of $\mathbb{B}$, and let $\tilde S_n$ denote a centered Gaussian element of $\mathbb{B}$ with the same covariance as $S_n$. Broadly speaking, the literature on rates of Gaussian approximation is concerned with bounding the  quantities 
 $\rho_n(\mathscr{A})=\sup_{A\in\mathscr{A}}\rho_n(A)$, or $\|\Delta_n\|_{\infty}=\sup_{t\in\R}|\Delta_n(t)|$, where
%
 \begin{align}
 \rho_n(A) & \ = \  \Big|\P(S_n\in A) - \P(\tilde S_n\in A)\Big|,\label{eqn:rhodef}\\[0.2cm]
 %
 \Delta_n(t) & \ = \ \Big|\P\big(\varphi(S_n)\leq t\big) - \P\big(\varphi(\tilde S_n)\leq t\big)\Big|.\label{eqn:deltadef}
 \end{align}
  In particular, note that if $\mathbb{B}=\R^p$ and $\tilde T=\max_{1\leq j\leq p} \tilde S_{n,j}$, then the distance $d_{\text{K}}(\mathcal{L}(T),\mathcal{L}(\tilde T))$ can be represented as either $\rho_n(\mathscr{A})$ or $\|\Delta_n\|_{\infty}$, by taking $\mathscr{A}$ to be a class of rectangles, or by taking $\varphi$ to be the coordinate-wise maximum function.
  %
  %
  
   Typically, the rates at which $\rho_n(\mathscr{A})$ and $\|\Delta_n\|_{\infty}$ decrease with $n$ is dependent on the distribution of $S_n$, the dimension of $\mathbb{B}$, as well as the smoothness of $\mathscr{A}$ and $\varphi$, among other factors. Although the study of $\rho_n(\mathscr{A})$ and $\|\Delta_n\|_{\infty}$ is highly multifaceted, it is a general principle that the smoothness of $\mathscr{A}$ and $\varphi$ tends to be much more influential when $\mathbb{B}$ is infinite-dimensional, as compared to the finite-dimensional case. (A discussion may be found in~\citep{Bentkus:1993}.) Indeed, this is worth emphasizing in relation to our work, since the choices of $\mathscr{A}$ and $\varphi$ corresponding to $d_{\text{K}}(\mathcal{L}(T),\mathcal{L}(\tilde T))$ are notable for their lack of smoothness.
   
   To illustrate this point, consider the following two facts that are known to hold when $\mathbb{B}$ is separable and infinite-dimensional~\citep{Bentkus:1993}: (1) If $\E[\|X_1\|_{\mathbb{B}}^3]<\infty$, and $\varphi$ satisfies smoothness conditions stronger than having three Fr\'echet derivatives,  then $\|\Delta_n\|_{\infty}$ is of order $\mathcal{O}(n^{-1/2})$. (2) If the previous conditions hold, except that $\varphi$ is only assumed to have one Fr\'echet derivative, then an example can be constructed so that $\|\Delta_n\|_{\infty}$ is bounded below by a sequence that converges to 0 arbitrarily slowly.  (Other examples of lower bounds may be found in the papers~\citep{Bentkus:1986,Talagrand:1984}, among others.) By contrast, in the finite-dimensional case where $\mathbb{B}=\R^p$ with $p$ held fixed as $n\to\infty$ and $\E[\|X_1\|_2^3]<\infty$, it is known that  non-smooth choices of $\varphi$ and $\mathscr{A}$ can lead to a $n^{-1/2}$ rate. Namely, it is known that $\|\Delta_n\|_{\infty}$ and $\rho_n(\mathscr{A})$ are of order $\mathcal{O}(n^{-1/2})$ when $\varphi$ is convex, or when $\mathscr{A}$ is the class of Borel convex sets. Beyond the case where $p$ is held fixed, many other works allow $n$ and $p$ to diverge together. When $p$ grows relatively slowly compared to $n$, the leading rate for the choices of $\varphi$ and $\mathscr{A}$ just mentioned is essentially $\mathcal{O}(p^{7/4}/n^{1/2})$~\citep{Bentkus:2003,Bentkus:2005}. See also~\citep{Sazonov:1968,Nagaev:1976,Senatov:1981,Portnoy:1986,Gotze:1991,Chen:2011,Zhai:2018} for additional background. Meanwhile, when $p\gg n$, it has recently been established that rates of the form $\mathcal{O}(\log(p)^b/n^{1/6})$ can be achieved if $\varphi$ is the coordinate-wise maximum function, or if $\mathscr{A}$ is a class of ``sparsely convex sets'', such as rectangles~\citep{CCK:AOP}.

   
   In light of the previous paragraph, it is natural that results in the infinite-dimensional setting have focused predominantly on smooth choices of $\mathscr{A}$ and $\varphi$. Nevertheless, there are some special cases where notable results have been obtained for non-smooth choices of $\varphi$ and $\mathscr{A}$ that correspond to max statistics. One such result is established in the paper~\citep{Asriev:1986}, which deals with bounds on $\rho_n(A(\mathbf{r}))$ for rectangular sets such as $A(\mathbf{r}):= \prod_{j=1}^{\infty}(-\infty,r_j]$,  with $\mathbf{r}=(r_1,r_2,\dots)$, in the infinite-dimensional Euclidean space $\R^{\infty}$. More specifically, if we continue to let $\sigma_j^2=\var(X_{1,j})$ for $j=1,2,\dots$, and define the ``effective dimension'' parameter
   $ {\tt{d}}(\mathbf{r})=\sum_{j=1}^{\infty} \ts\frac{\sigma_j}{\sigma_j+r_j},$
   then the following bound holds under certain conditions on the distribution of $S_n$,
   \begin{equation}\label{eqn:asriev}
    \rho_n(A(\mathbf{r})) \ \lesssim \  \ts\frac{\log(n)^{3/2}}{n^{1/2}} \cdot {\tt{d}}(\mathbf{r})\cdot(1+\tt{d}(\mathbf{r})^2).
   \end{equation}
To comment on how the bound~\eqref{eqn:asriev} relates with our Gaussian approximation result in Theorem~3.1, note that they both involve near-parametric rates, and are governed by the parameters $(\sigma_1,\sigma_2,\dots)$. However, there are also some crucial differences. First, the bound~\eqref{eqn:asriev} is \emph{non-uniform} with respect to the set $A(\mathbf{r})$, whereas Theorem~3.1 is a uniform result. In particular, this difference becomes apparent by setting all $r_j$ equal to $\sigma_{(n)}$, which implies ${\tt{d}}(\mathbf{r})\geq n/2$, and causes the bound~\eqref{eqn:asriev} diverge as $n\to\infty$.  Second, the bound~\eqref{eqn:asriev} relies on the assumption that $S_n$ has a \emph{diagonal covariance} matrix, whereas Theorem~3.1 allows for much more general covariance structures.

Some further examples related to max statistics arise in the context of empirical process theory. Let $\G_n(f)=\sqrt{n}\sum_{i=1}^n (f(X_i)-\E[f(X_i)])$ denote an empirical process that is generated from i.i.d.~observations $X_1,\dots,X_n$, and indexed by a class of functions $f\in\mathscr{F}$. Also let the Gaussian counterpart of $\G_n$ be denoted by $\tilde{\G}$ (i.e.~a Brownian bridge)~\citep[][]{vanderVaart:Wellner:2000}. In this setting, the paper~\citep{Norvaivsa:1991} studies the quantity
\begin{equation}\label{eqn:Deltanew}
\Delta_n(t)=\Big|\P\big(\|\G_n\|_{\mathscr{F}}\leq t\big) - \P\big(\|\tilde\G\|_{\mathscr{F}}\leq t\big)\Big|,
\end{equation}
which can be understood in terms of the earlier definition~\eqref{eqn:deltadef} by letting $\varphi(\G_n)=\|\G_n\|_{\mathscr{F}}=\sup_{f\in\mathscr{F}}|\G_n(f)|$. Under the assumption that $\mathscr{F}$ is a VC subgraph class of uniformly bounded functions, it is shown in the paper~\citep{Norvaivsa:1991} that the bound
\begin{equation}\label{eqn:empbound1}
\Delta_n(t) \ \lesssim \ \ts\frac{1}{(1+t)^3}\ts\frac{\log(n)^2}{n^{1/6}}
\end{equation}
holds for all $t\geq 0$. In relation to the current work, the condition that the functions in $\mathscr{F}$ are uniformly bounded is an important distinction, because the max statistic $T=\max_{1\leq j\leq p} \G_n(f_j)$ arises from the functions $f_j(x)=x_j$, which are unbounded. 

 More recently, the papers~\citep{CCK:suprema,CCK:SPA} have derived bounds on the quantity~$\|\Delta_n\|_{\infty}$ (as well as coupling probabilities) under weaker assumptions on the class $\mathscr{F}$. For instance, these works allow for classes of functions that are non-Donsker or unbounded, provided that a suitable envelope function is available. Nevertheless, the resulting bounds on $\|\Delta_n\|_{\infty}$ involve restrictions on how quickly the parameter $\inf_{f\in\mathscr{F}}\var(\G_n(f))$ can decrease with $n$ --- whereas our results do not involve such restrictions. Also, the rates developed in these works are broadly similar to~\eqref{eqn:empbound1}, and it is unclear if a modification of the techniques would lead to near-parametric rates in our setting.
Lastly, a number of earlier results that are related to bounding $\|\Delta_n\|_{\infty}$, such as invariance principles and couplings, can be found in the papers~\citep[][and references therein]{Csorgo:1986,Massart:1986,Massart:1989,Paulauskas:Stieve:1990,Bloznelis:1997}. However, these works are tailored to fairly specialized processes (e.g., dealing with c\`adl\`ag functions on the unit interval, rectangles in $\R^p$ when $p$ is fixed, or processes generated by Uniform[0,1] random variables).






\end{document}
